\documentclass[10pt]{article}
\usepackage[utf8]{inputenc}
\usepackage[T1]{fontenc}
\usepackage{amsmath}
\usepackage{amsfonts}
\usepackage{amssymb}
\usepackage[version=4]{mhchem}
\usepackage{stmaryrd}
\usepackage{hyperref}
\hypersetup{colorlinks=true, linkcolor=blue, filecolor=magenta, urlcolor=cyan,}
\urlstyle{same}
\usepackage{bbold}
\usepackage{caption}

\title{New Classes of Almost Bent and Almost Perfect Nonlinear Polynomials* }

\author{Lilya Budaghyan, Claude Carlet ${ }^{\ddagger}$, Alexander Pott ${ }^{\S}$}
\date{}


%New command to display footnote whose markers will always be hidden
\let\svthefootnote\thefootnote
\newcommand\blfootnotetext[1]{%
  \let\thefootnote\relax\footnote{#1}%
  \addtocounter{footnote}{-1}%
  \let\thefootnote\svthefootnote%
}

%Overriding the \footnotetext command to hide the marker if its value is `0`
\let\svfootnotetext\footnotetext
\renewcommand\footnotetext[2][?]{%
  \if\relax#1\relax%
    \ifnum\value{footnote}=0\blfootnotetext{#2}\else\svfootnotetext{#2}\fi%
  \else%
    \if?#1\ifnum\value{footnote}=0\blfootnotetext{#2}\else\svfootnotetext{#2}\fi%
    \else\svfootnotetext[#1]{#2}\fi%
  \fi
}

\begin{document}
\maketitle
\captionsetup{singlelinecheck=false}


\begin{abstract}
We construct infinite classes of almost bent and almost perfect nonlinear polynomials, which are affinely inequivalent to any sum of a power function and an affine function.
\end{abstract}

Keywords: Vectorial Boolean function, Nonlinearity, Differential uniformity, Almost perfect nonlinear, Almost bent, Affine equivalence, CCZ-equivalence.

\section*{1 Introduction}
Vectorial Boolean functions are used in cryptography, more precisely in block ciphers. An important condition on these functions is a high resistance to the differential and linear cryptanalyses, which are the main attacks on block ciphers. The functions with the smallest possible differential uniformity (see [26]) oppose an optimum resistance to the differential attack 3]. They are called almost perfect nonlinear (APN). The functions achieving the maximal possible nonlinearity (see [13, 25]) possess the best resistance to the linear attack [24] and they are called almost bent (AB) or maximum nonlinear.

Up to now, in the study of APN and AB functions the main attention has been payed to power mappings and all known constructions of APN and AB functions happen to be equivalent to power functions (see for example [5, 6, 7, 16, 17, 18]). Recall that functions $F$ and $F^{\prime}$ are called equivalent if either $F^{\prime}=A_{1} \circ F \circ A_{2}+A$ or $F^{\prime}=A_{1} \circ F^{-1} \circ A_{2}+A$ (in case $F$ is a permutation) for some affine functions $A, A_{1}$ and $A_{2}$, where $A_{1}$ and $A_{2}$ are permutations. We shall say that two functions $F$ and $F^{\prime}=A_{1} \circ F \circ A_{2}+A$ (with $A, A_{1}, A_{2}$ affine, $A_{1}, A_{2}$ being permutations) are extended affine equivalent (EA-equivalent). In this

\footnotetext{*Part of this paper was presented at WCC 2005 [4].\\
${ }^{\dagger}$ Institute of Algebra and Geometry, Otto-von-Guericke University Magdeburg, D-39016 Magdeburg, GERMANY; e-mail: \href{mailto:lilya.budaghyan@student.uni-magdeburg.de}{lilya.budaghyan@student.uni-magdeburg.de}.\\
${ }^{\ddagger}$ INRIA, Projet CODES, BP 105-78153, Le Chesnay Cedex, FRANCE; e-mail: \href{mailto:claude.carlet@inria.fr}{claude.carlet@inria.fr}; also member of the University of Paris 8.\\
${ }^{§}$ Institute of Algebra and Geometry, Otto-von-Guericke University Magdeburg, D-39016 Magdeburg, GERMANY; e-mail: \href{mailto:alexander.pott@mathematik.uni-magdeburg.de}{alexander.pott@mathematik.uni-magdeburg.de}.
}
paper we give the first theoretical constructions of APN and AB polynomials which are inequivalent to power mappings. In these constructions, we apply the transformation of functions given by Carlet, Charpin and Zinoviev (see Proposition 3 in [11]) to the Gold APN and AB mappings [19, 26]. This transformation leads to an equivalence relation of functions that we call the Carlet-Charpin-Zinoviev equivalence (CCZ-equivalence). CCZequivalence corresponds to the affine equivalence of the graphs of functions, i.e. functions $F$ and $F^{\prime}$ are CCZ-equivalent if and only if, for some affine permutation, the image of the graph of $F$ is the graph of the function $F^{\prime}$. CCZ-equivalence preserves the nonlinearity, the differential uniformity of functions and the resistance of a function to the algebraic cryptanalysis [14. It could be expected that CCZ-equivalence coincides in practice with the equivalence mentioned above. The present paper aims at proving and illustrating that CCZ-equivalence is more general.

The next section contains all the necessary definitions related to vectorial Boolean functions, including EA-equivalence, APN and AB properties.

In Section 3, we give the definition of CCZ-equivalence, we describe its main properties and we show its connections with EA-equivalence.

We give some results related to a classification of functions CCZ-equivalent to the Gold mappings in Section 4.

Theorems 1, 2 and 4 in Section 5 present constructions of AB and APN polynomials which are EA-inequivalent to power functions and Theorem 3 presents an APN function EA-inequivalent to all known APN mappings.

This paper is an extended version of an abstract published in [4]. In particular, Theorem 4, the proof of Theorem 3 and some properties of CCZ-equivalence are not included in the abstract.

\section*{2 Almost perfect nonlinear and almost bent functions}
Let $\mathbb{F}_{2}^{m}$ be the $m$-dimensional vector space over the field $\mathbb{F}_{2}$. Any function $F$ from $\mathbb{F}_{2}^{m}$ to itself can be uniquely represented as a polynomial on $m$ variables with coefficients in $\mathbb{F}_{2}^{m}$, whose degree with respect to each coordinate is at most 1 :

$$
F\left(x_{1}, \ldots, x_{m}\right)=\sum_{u \in \mathbb{F}_{2}^{m}} c(u)\left(\prod_{i=1}^{m} x_{i}^{u_{i}}\right), \quad c(u) \in \mathbb{F}_{2}^{m}
$$

This representation is called the algebraic normal form of $F$ and its degree $d^{\circ}(F)$ the algebraic degree of the function $F$.\\
Besides, viewed as a function from the field $\mathbb{F}_{2^{m}}$ to itself, $F$ has a unique representation as a univariate polynomial over $\mathbb{F}_{2^{m}}$ of degree smaller than $2^{m}$ :

$$
F(x)=\sum_{i=0}^{2^{m}-1} c_{i} x^{i}, \quad c_{i} \in \mathbb{F}_{2^{m}}
$$

For any $k, 0 \leq k \leq 2^{m}-1$, the number $w_{2}(k)$ of the nonzero coefficients $k_{s} \in\{0,1\}$ in the binary expansion $\sum_{s=0}^{m-1} 2^{s} k_{s}$ of $k$ is called the 2 -weight of $k$. The algebraic degree of $F$ is equal to the maximum 2 -weight of the exponents $i$ of the polynomial $F(x)$ such that $c_{i} \neq 0$, that is $d^{\circ}(F)=\max _{0 \leq i \leq n-1, c_{i} \neq 0} w_{2}(i)$ (see [11]).

A function $F: \mathbb{F}_{2}^{m} \rightarrow \mathbb{F}_{2}^{m}$ is linear if and only if $F(x)$ is a linearized polynomial over $\mathbb{F}_{2^{m}}$, that is,

$$
\sum_{i=0}^{m-1} c_{i} x^{2^{i}}, \quad c_{i} \in \mathbb{F}_{2^{m}}
$$

The sum of a linear function and a constant is called an affine function.\\
Let $F$ be a function from $\mathbb{F}_{2}^{m}$ to itself and $A_{1}, A_{2}: \mathbb{F}_{2}^{m} \rightarrow \mathbb{F}_{2}^{m}$ be affine permutations. Then the functions $F$ and $A_{1} \circ F \circ A_{2}$ are called affine equivalent. Affine equivalent functions have the same algebraic degree (i.e. the algebraic degree is affine invariant).

We shall say that the functions $F$ and $F^{\prime}$ are extended affine equivalent (EA-equivalent) if $F^{\prime}=A_{1} \circ F \circ A_{2}+A$ for some affine permutations $A_{1}, A_{2}$ and an affine function $A$. If $F$ is not affine, then $F$ and $F^{\prime}$ have again the same algebraic degree.

For a function $F: \mathbb{F}_{2}^{m} \rightarrow \mathbb{F}_{2}^{m}$ and any elements $a, b \in \mathbb{F}_{2}^{m}$ we denote

$$
\delta_{F}(a, b)=\left|\left\{x \in \mathbb{F}_{2}^{m}: F(x+a)+F(x)=b\right\}\right|
$$

and

$$
\Delta_{F}=\left\{\delta_{F}(a, b): a, b \in \mathbb{F}_{2}^{m}, a \neq 0\right\}
$$

$F$ is called a differentially $\delta$-uniform function if $\max _{a \in \mathbb{F}_{2}^{m *}, b \in \mathbb{F}_{2}^{m}} \delta_{F}(a, b) \leq \delta$, where $\mathbb{F}_{2}^{m *}= \mathbb{F}_{2}^{m} \backslash\{0\}$. For any $a, b \in \mathbb{F}_{2}^{m}$, the number $\delta_{F}(a, b)$ is even since if $x_{0}$ is a solution of the equation $F(x+a)+F(x)=b$ then $x_{0}+a$ is a solution too. Hence, $\delta \geq 2$. Differentially 2 -uniform mappings are called almost perfect nonlinear.

For any function $F: \mathbb{F}_{2}^{m} \rightarrow \mathbb{F}_{2}^{m}$, the values

$$
\lambda_{F}(a, b)=\sum_{x \in \mathbb{F}_{2}^{m}}(-1)^{b \cdot F(x)+a \cdot x}, \quad a, b \in \mathbb{F}_{2}^{m}
$$

do not depend on a particular choice of the inner product $" . "$ in $\mathbb{F}_{2}^{m}$. If we identify $\mathbb{F}_{2}^{m}$ with $\mathbb{F}_{2^{m}}$ then we can take $x \cdot y=\operatorname{tr}(x y)$, where $\operatorname{tr}(x)=x+x^{2}+x^{4}+\ldots+x^{2^{m-1}}$ is the trace function from $\mathbb{F}_{2^{m}}$ into $\mathbb{F}_{2}$. Later we shall also need for a divisor $n$ of $m$ the relative trace $\operatorname{tr}_{m / n}(x)=x+x^{2^{n}}+x^{2^{2 n}}+\ldots+x^{2^{(m / n-1) n}}$ and we denote $\operatorname{tr}_{n}(x)=x+x^{2}+x^{2^{2}}+\ldots+x^{2^{n-1}}$. The set $\Lambda_{F}=\left\{\lambda_{F}(a, b): a, b \in \mathbb{F}_{2}^{m}, b \neq 0\right\}$ is called the Walsh spectrum of $F$ and the value

$$
\mathcal{N} \mathcal{L}(F)=2^{m-1}-\frac{1}{2} \max _{a \in \mathbb{F}_{2}^{m}, b \in \mathbb{F}_{2}^{m *}}\left|\lambda_{F}(a, b)\right|
$$

the nonlinearity of the function $F$. The nonlinearity of any function $F$ satisfies the inequality

$$
\mathcal{N} \mathcal{L}(F) \leq 2^{m-1}-2^{\frac{m-1}{2}}
$$

(13, 29) and in case of equality $F$ is called almost bent or maximum nonlinear. For any AB function $F$, the Walsh spectrum $\Lambda_{F}$ equals $\left\{0, \pm 2^{\frac{m+1}{2}}\right\}$ as it is proven in 13 .

For EA-equivalent functions $F$ and $F^{\prime}$, we have $\Delta_{F}=\Delta_{F^{\prime}}, \Lambda_{F}=\Lambda_{F^{\prime}}$ and if $F$ is a permutation then $\Delta_{F}=\Delta_{F^{-1}}, \Lambda_{F}=\Lambda_{F^{-1}}$ (see [11]). Therefore, if $F$ is APN (resp. AB) and $F^{\prime}$ is EA-equivalent to either $F$ or $F^{-1}$ (if $F$ is a permutation), then $F^{\prime}$ is also APN (resp. AB).

Table 1 (resp. Table 2) gives all known values of exponents $d$ (up to EA-equivalence and up to taking the inverse when a function is a permutation) such that the power function $x^{d}$ is APN (resp. AB).

\begin{table}[h]
\begin{center}
\captionsetup{labelformat=empty}
\caption{Table 1\\
Known APN power functions $x^{d}$ on $\mathbb{F}_{2^{m}}$.}
\begin{tabular}{|l|l|l|l|l|}
\hline
 & Exponents $d$ & Conditions & $w_{2}(d)$ & Proven in \\
\hline
Gold functions & $2^{i}+1$ & $\operatorname{gcd}(i, m)=1$ & 2 & 19, 26 \\
\hline
Kasami functions & $2^{2 i}-2^{i}+1$ & $\operatorname{gcd}(i, m)=1$ & $i+1$ & 21, 22] \\
\hline
Welch function & $2^{t}+3$ & $m=2 t+1$ & 3 & [17] \\
\hline
Niho function & \( \begin{gathered} 2^{t}+2^{\frac{t}{2}}-1, t \text { even } \\ 2^{t}+2^{\frac{3 t+1}{2}}-1, t \text { odd } \end{gathered} \) & $m=2 t+1$ & \( \begin{gathered} (t+2) / 2 \\ t+1 \end{gathered} \) & 16 \\
\hline
Inverse function & $2^{2 t}-1$ & $m=2 t+1$ & $m-1$ & 226 \\
\hline
Dobbertin function & $2^{4 i}+2^{3 i}+2^{2 i}+2^{i}-1$ & $m=5 i$ & $i+3$ & [18] \\
\hline
\end{tabular}
\end{center}
\end{table}

\begin{table}[h]
\begin{center}
\captionsetup{labelformat=empty}
\caption{Table 2\\
Known AB power functions $x^{d}$ on $\mathbb{F}_{2^{m}}, m$ odd.}
\begin{tabular}{|c|c|c|c|}
\hline
 & Exponents $d$ & Conditions & Proven in \\
\hline\hline
Gold functions & $2^{i}+1$ & $\operatorname{gcd}(i, m)=1$ & $[19,[26]$ \\
\hline
Kasami functions & $2^{2 i}-2^{i}+1$ & $\operatorname{gcd}(i, m)=1$ & $[22]$ \\
\hline
Welch function & $2^{t}+3$ & $m=2 t+1$ & $[6,[7]$ \\
\hline
Niho function & \begin{tabular}{l}
$2^{t}+2^{\frac{t}{2}}-1, t$ even \\
$2^{t}+2^{\frac{3 t+1}{2}}-1, t$ odd \\
\end{tabular} & $m=2 t+1$ & $[20]$ \\
\hline
\end{tabular}
\end{center}
\end{table}

Every AB function is APN [13. The converse is not true since AB functions exist only when $m$ is odd while APN functions exist for $m$ even too. Besides, in the $m$ odd case, the Dobbertin APN function is not AB as proven in [7]. Also the inverse APN function is not AB since it has the algebraic degree $m-1$ while the algebraic degree of any AB function is not greater than $(m+1) / 2$ (see [11]). When $m$ is even, the inverse function $x^{2^{m}-2}$ is a differentially 4 -uniform permutation 26 and has the best known nonlinearity 23, that is $2^{m-1}-2^{\frac{m}{2}}$ (see [7, 15]). When $m \equiv 2[\bmod 4]$ and $\operatorname{gcd}(i, m)=2$, the functions $x^{2^{i}+1}$ and $x^{2^{2 i}-2^{i}+1}$ also have the best known nonlinearity 19, 22 and when $\operatorname{gcd}(i, m)=s$ and $m / s$ is odd, these functions have three valued Walsh spectrum $\left\{0, \pm 2^{\frac{m+s}{2}}\right\}$ (26, 30).

\section*{3 Carlet-Charpin-Zinoviev equivalence of functions}
For a function $F$ from $\mathbb{F}_{2}^{m}$ to itself, we denote by $G_{F}$ the graph of the function $F$ :

$$
G_{F}=\left\{(x, F(x)): x \in \mathbb{F}_{2}^{m}\right\} \subset \mathbb{F}_{2}^{2 m} .
$$

The property of stability of APN and AB functions given in Proposition 3 of 11 leads to the definition of the following equivalence relation of functions.

Definition 1 (CCZ-equivalence) We say that functions $F, F^{\prime}: \mathbb{F}_{2}^{m} \rightarrow \mathbb{F}_{2}^{m}$ are Carlet-Charpin-Zinoviev equivalent if there exists a linear (affine) permutation $\mathcal{L}: \mathbb{F}_{2}^{2 m} \rightarrow \mathbb{F}_{2}^{2 m}$ such that $\mathcal{L}\left(G_{F}\right)=G_{F^{\prime}}$.

We shall consider only the case of linear functions, but all proofs for the statements related to the CCZ-equivalence can be easily extended for the case of affine functions too.

A linear function $\mathcal{L}: \mathbb{F}_{2}^{2 m} \rightarrow \mathbb{F}_{2}^{2 m}$ can be considered as a pair of linear functions $L_{1}, L_{2}: \mathbb{F}_{2}^{2 m} \rightarrow \mathbb{F}_{2}^{m}$ such that $\mathcal{L}(x, y)=\left(L_{1}(x, y), L_{2}(x, y)\right)$ for all $x, y \in \mathbb{F}_{2}^{m}$. Then for a function $F: \mathbb{F}_{2}^{m} \rightarrow \mathbb{F}_{2}^{m}$ we have $\mathcal{L}(x, F(x))=\left(F_{1}(x), F_{2}(x)\right)$, where


\begin{align*}
& F_{1}(x)=L_{1}(x, F(x))  \tag{1}\\
& F_{2}(x)=L_{2}(x, F(x)) \tag{2}
\end{align*}


Hence, the set $\mathcal{L}\left(G_{F}\right)=\left\{\left(F_{1}(x), F_{2}(x)\right): x \in \mathbb{F}_{2}^{m}\right\}$ is the graph of a function $F^{\prime}$ if and only if the function $F_{1}$ is a permutation. If $\mathcal{L}$ and $F_{1}$ are permutations then $\mathcal{L}\left(G_{F}\right)=G_{F^{\prime}}$, where $F^{\prime}=F_{2} \circ F_{1}^{-1}$, and the functions $F$ and $F^{\prime}$ are CCZ-equivalent.

In the proposition below, we give a slightly different approach to the CCZ-equivalence, that we shall use in the constructions of APN polynomials which will be EA-inequivalent to power functions. Recall that a set $G \subset \mathbb{F}_{2}^{2 m}$ is transversal to a subgroup $V$ of $\mathbb{F}_{2}^{2 m}$ if $|G \cap(u+V)|=1$ for any $u \in \mathbb{F}_{2}^{2 m}$.

Proposition 1 Let $F: \mathbb{F}_{2}^{m} \rightarrow \mathbb{F}_{2}^{m}$ and $G \subset \mathbb{F}_{2}^{2 m}$. Then the set $G$ is the graph of a function CCZ-equivalent to $F$ if and only if there exists a linear permutation $\mathcal{L}: \mathbb{F}_{2}^{2 m} \rightarrow \mathbb{F}_{2}^{2 m}$ such that $G=\mathcal{L}\left(G_{F}\right)$ and $G_{F}$ is transversal to $V^{\prime}=\mathcal{L}^{-1}(V)$, where $V=\left\{(0, x): x \in \mathbb{F}_{2}^{m}\right\}$.

Proof. The condition that there exists a linear permutation $\mathcal{L}: \mathbb{F}_{2}^{2 m} \rightarrow \mathbb{F}_{2}^{2 m}$ such that $G=\mathcal{L}\left(G_{F}\right)$ is clearly necessary. Let us denote $U=\left\{(x, 0): x \in \mathbb{F}_{2}^{m}\right\}, V=\{(0, x): x \in \left.\mathbb{F}_{2}^{m}\right\}, \mathcal{L}^{-1}(U)=U^{\prime}$ and $\mathcal{L}^{-1}(V)=V^{\prime}$. Then $G$ is the graph of a function if and only if $|G \cap(u+V)|=1$ for any $u \in U$; that is, if and only if $\left|\mathcal{L}^{-1}(G) \cap\left(u^{\prime}+V^{\prime}\right)\right|=1$ for any $u^{\prime} \in U^{\prime}$. Hence, $G$ is the graph of a function CCZ-equivalent to $F$ if and only if $G_{F}$ is transversal to $V^{\prime}$.

Let $F$ be an arbitrary function on $\mathbb{F}_{2}^{m}$ and $V^{\prime}$ be an arbitrary subgroup of $\mathbb{F}_{2}^{2 m}$. If we denote

$$
\begin{aligned}
& H_{a}=\left\{F(x+a)+F(x): x \in \mathbb{F}_{2}^{m}\right\}, \quad a \in \mathbb{F}_{2}^{m *} \\
& A_{F}=\left\{(a, F(x+a)+F(x)): a \in \mathbb{F}_{2^{m}}^{*}, x \in \mathbb{F}_{2^{m}}\right\}
\end{aligned}
$$

then $A_{F}=\bigcup_{a \in \mathbb{F}_{2}^{*} m}\left(a, H_{a}\right)$. It is easy to note that $A_{F} \cap V^{\prime}=\emptyset$ if and only if $G_{F}$ is transversal to $V^{\prime}$. We shall see that, to construct a function CCZ-equivalent but EA-inequivalent to\\
the function $F$, it is necessary that the subgroup $V^{\prime}$ in Proposition 1 is different from $V$ and as soon as we have a subroup $V^{\prime}$ such that $G_{F}$ is transversal to $V^{\prime}$ we can find a linear permutation $\mathcal{L}$ that $V^{\prime}=\mathcal{L}^{-1}(V)$. However, even if we have such a subgroup it is difficult to determine whether the resulting function is EA-inequivalent to $F$.

The property of stability of APN and AB mappings (see 11) can be easily generalized to all functions (not necessarily APN or AB) as follows:

Proposition 2 Let $F, F^{\prime}: \mathbb{F}_{2}^{m} \rightarrow \mathbb{F}_{2}^{m}$ be $C C Z$-equivalent functions. Then $\Delta_{F}=\Delta_{F^{\prime}}$ and $\Lambda_{F}=\Lambda_{F^{\prime}}$. In particular, $F$ is an APN (resp. AB) function if and only if $F^{\prime}$ is APN (resp. $A B$ ).

Proof. If $F, F^{\prime}: \mathbb{F}_{2}^{m} \rightarrow \mathbb{F}_{2}^{m}$ are CCZ-equivalent, then $F^{\prime}=F_{2} \circ F_{1}^{-1}$ for a certain linear permutation $\mathcal{L}=\left(L_{1}, L_{2}\right)$, where $F_{1}, F_{2}$ are defined by (11) and (2). For any $(a, b) \in \mathbb{F}_{2}^{2 m}$ we have

$$
\begin{aligned}
\lambda_{F}(a, b) & =\sum_{x \in \mathbb{F}_{2}^{m}}(-1)^{b \cdot F(x)+a \cdot x}=\sum_{x \in \mathbb{F}_{2}^{m}}(-1)^{(a, b) \cdot(x, F(x))} \\
& =\sum_{x \in \mathbb{F}_{2}^{m}}(-1)^{(a, b) \cdot \mathcal{L}^{-1}\left(F_{1}(x), F_{2}(x)\right)}=\sum_{x \in \mathbb{F}_{2}^{m}}(-1)^{\mathcal{L}^{-1 *}(a, b) \cdot\left(x, F_{2} \circ F_{1}^{-1}(x)\right)} \\
& =\lambda_{F^{\prime}}\left(\mathcal{L}^{-1 *}(a, b)\right),
\end{aligned}
$$

where $\mathcal{L}^{-1 *}$ is the adjoint operator of $\mathcal{L}^{-1}$ (i.e. $x \cdot \mathcal{L}^{-1}(y)=\mathcal{L}^{-1 *}(x) \cdot y$, for any $(x, y) \in \mathbb{F}_{2}^{2 m}$; if "." is the usual inner product, then $\mathcal{L}^{-1 *}$ is the linear permutation whose matrix is transposed of that of $\mathcal{L}^{-1}$ ). Hence, $\Lambda_{F}=\Lambda_{F^{\prime}}$.\\
The proof that $\Delta_{F}=\Delta_{F^{\prime}}$ for arbitrary functions $F, F^{\prime}$ is completely the same as in the case when one of the functions $F, F^{\prime}$ is APN (see [11]).

Remark 1 Obviously, CCZ-equivalence can be defined for functions between any two groups $H_{1}$ and $H_{2}$. For a function $F: H_{1} \rightarrow H_{2}$ we can consider the set of the values $\delta_{F}(a, b)=\left|\left\{x \in H_{1}: F(x+a)-F(x)=b\right\}\right|, a \in H_{1} \backslash\{0\}, b \in H_{2}$, and if the groups $H_{1}$ and $H_{2}$ are Abelian, then the discrete Fourier transform of $F$ can also be defined. In this more general case CCZ-equivalent functions again have the same differential and linear properties. One can find results related to nonlinear functions on Abelian groups in [12, 28].

Since CCZ-equivalent functions have the same differential uniformity and the same nonlinearity, the resistance of a function to linear and differential attacks is CCZ-invariant. CCZ-equivalent functions have also the same weakness/strength with respect to algebraic attacks. Indeed, let functions $F, F^{\prime}: \mathbb{F}_{2}^{m} \rightarrow \mathbb{F}_{2}^{m}$ be CCZ-equivalent. Then $F^{\prime}=F_{2} \circ F_{1}^{-1}$, where $F_{1}, F_{2}$ are defined by (1) and (2) for a certain linear permutation $\mathcal{L}=\left(L_{1}, L_{2}\right)$. If there exists a nonzero function $\psi: \mathbb{F}_{2}^{2 m} \rightarrow \mathbb{F}_{2}$ of low degree such that


\begin{equation*}
\psi(x, F(x))=0, \quad \forall x \in \mathbb{F}_{2}^{m} \tag{3}
\end{equation*}


then $\psi$ could be used in an algebraic attack [14]. The function $\psi \circ \mathcal{L}^{-1}$ has the same degree as $\psi$ and Relation (3) is equivalent to

$$
\psi \circ \mathcal{L}^{-1}\left(F_{1}(x), F_{2}(x)\right)=0, \quad \forall x \in \mathbb{F}_{2}^{m}
$$

which implies

$$
\psi \circ \mathcal{L}^{-1}\left(x, F^{\prime}(x)\right)=0, \quad \forall x \in \mathbb{F}_{2}^{m}
$$

by replacing $x$ by $F_{1}^{-1}(x)$. Hence, $\psi \circ \mathcal{L}^{-1}$ could be used in an algebraic attack on $F^{\prime}$, and vice versa. Therefore, the resistance of a function to algebraic attacks is also CCZinvariant.\\
Since any permutation is CCZ-equivalent to its inverse then obviously the algebraic degree of a function is not CCZ-invariant (while it is EA-invariant as noted above). For example, if $F: \mathbb{F}_{2}^{m} \rightarrow \mathbb{F}_{2}^{m}$ is a Gold AB function then $d^{\circ}(F)=2$ and $d^{\circ}\left(F^{-1}\right)=\frac{m+1}{2}$ as proven in 26.

EA-equivalent functions are CCZ-equivalent and if a function $F$ is a permutation then $F$ is CCZ-equivalent to $F^{-1}$ 11. More precisely:

Proposition 3 Let $F, F^{\prime}: \mathbb{F}_{2}^{m} \rightarrow \mathbb{F}_{2}^{m}$. The function $F^{\prime}$ is $E A$-equivalent to the function $F$ or to the inverse of $F$ (if it exists) if and only if there exists a linear permutation $\mathcal{L}=\left(L_{1}, L_{2}\right)$ on $\mathbb{F}_{2}^{2 m}$ such that $\mathcal{L}\left(G_{F}\right)=G_{F^{\prime}}$ and the function $L_{1}$ depends only on one variable, i.e. $L_{1}(x, y)=L(x)$ or $L_{1}(x, y)=L(y)$.

Proof. Let $\mathcal{L}\left(G_{F}\right)=G_{F^{\prime}}$ for some linear permutation $\mathcal{L}=\left(L_{1}, L_{2}\right): \mathbb{F}_{2}^{2 m} \rightarrow \mathbb{F}_{2}^{2 m}$ and $L_{1}(x, y)=L(y), L_{2}(x, y)=R_{1}(x)+R_{2}(y)$, where $L, R_{1}, R_{2}: \mathbb{F}_{2}^{m} \rightarrow \mathbb{F}_{2}^{m}$ are linear. Then

$$
\begin{gathered}
F_{1}(x)=L_{1}(x, F(x))=L \circ F(x), \\
F_{2}(x)=L_{2}(x, F(x))=R_{1}(x)+R_{2} \circ F(x) .
\end{gathered}
$$

The function $F_{1}$ is a permutation, since $\mathcal{L}\left(G_{F}\right)$ is the graph of a function. Therefore, $L$ and $F$ have to be permutations. On the other hand, the system

$$
\begin{cases}L(y) & =0 \\ R_{1}(x)+R_{2}(y) & =0\end{cases}
$$

has only $(0,0)$ solution, since $\mathcal{L}=\left(L_{1}, L_{2}\right)$ is a permutation. But $L$ is also a permutation. Therefore, the linear function $R_{1}$ has to be a permutation too.\\
We have

$$
\begin{aligned}
F^{\prime}(x) & =F_{2} \circ F_{1}^{-1}(x)=R_{1} \circ F^{-1} \circ L^{-1}(x)+R_{2} \circ F \circ F^{-1} \circ L^{-1}(x) \\
& =R_{1} \circ F^{-1} \circ L^{-1}(x)+R_{2} \circ L^{-1}(x)
\end{aligned}
$$

Thus, $F^{\prime}$ is EA-equivalent to $F^{-1}$.\\
The proof that $F^{\prime}$ is EA-equivalent to $F$, when $L_{1}(x, y)=L(x)$, is similar.\\
Conversely, let $F^{\prime}=R_{1} \circ F \circ R_{2}+R$ or $F^{\prime}=R_{1} \circ F^{-1} \circ R_{2}+R$ for some linear permutations $R_{1}, R_{2}$ and for some linear function $R$. Then take $\mathcal{L}(x, y)=\left(R_{2}^{-1}(x), R_{1}(y)+R \circ R_{2}^{-1}(x)\right)$\\
in the first case and in the second case take $\mathcal{L}(x, y)=\left(R_{2}^{-1}(y), R \circ R_{2}^{-1}(y)+R_{1}(x)\right)$. Obviously, all conditions are satisfied.

Remark 2 Proposition 3 shows that, for a function $F: \mathbb{F}_{2}^{m} \rightarrow \mathbb{F}_{2}^{m}$, if $\mathcal{L}=\left(L_{1}, L_{2}\right)$ and $\mathcal{L}^{\prime}=\left(L_{1}, L_{2}^{\prime}\right)$ are linear permutations on $\mathbb{F}_{2}^{2 m}$ such that the function $L_{1}(x, F(x))$ is a permutation, then the functions defined by the graphs $\mathcal{L}\left(G_{F}\right)$ and $\mathcal{L}^{\prime}\left(G_{F}\right)$ are EAequivalent. Indeed, we have $\left(\mathcal{L}^{\prime} \circ \mathcal{L}^{-1}\right) \circ \mathcal{L}\left(G_{F}\right)=\mathcal{L}^{\prime}\left(G_{F}\right)$ and, denoting $\left(\mathcal{L}^{\prime} \circ \mathcal{L}^{-1}\right)=\left(S_{1}, S_{2}\right)$, we have $S_{1}(x, y)=x$. This last equality can be easily checked by considering the linear functions $\mathcal{L}, \mathcal{L}^{\prime}$ and $\mathcal{L}^{-1}$ as $(2 m) \times(2 m)$ matrices

$$
\mathcal{L}=\left(\begin{array}{ll}
R_{1} & R_{2} \\
T_{1} & T_{2}
\end{array}\right), \quad \mathcal{L}^{\prime}=\left(\begin{array}{ll}
R_{1} & R_{2} \\
T_{1}^{\prime} & T_{2}^{\prime}
\end{array}\right), \quad \mathcal{L}^{-1}=\left(\begin{array}{ll}
A_{1} & A_{2} \\
A_{3} & A_{4}
\end{array}\right) .
$$

We have

$$
\left(\begin{array}{ll}
R_{1} & R_{2} \\
T_{1} & T_{2}
\end{array}\right) \times\left(\begin{array}{ll}
A_{1} & A_{2} \\
A_{3} & A_{4}
\end{array}\right)=\left(\begin{array}{ll}
R_{1} A_{1}+R_{2} A_{3} & R_{1} A_{2}+R_{2} A_{4} \\
T_{1} A_{1}+T_{2} A_{3} & T_{1} A_{2}+T_{2} A_{4}
\end{array}\right)=\left(\begin{array}{ll}
I & 0 \\
0 & I
\end{array}\right),
$$

where $I$ is the identity matrix and 0 is the 0 -matrix of order $m$, and this implies

$$
\left(\begin{array}{ll}
R_{1} & R_{2} \\
T_{1}^{\prime} & T_{2}^{\prime}
\end{array}\right) \times\left(\begin{array}{ll}
A_{1} & A_{2} \\
A_{3} & A_{4}
\end{array}\right)=\left(\begin{array}{cc}
I & 0 \\
T_{1}^{\prime} A_{1}+T_{2}^{\prime} A_{3} & T_{1}^{\prime} A_{2}+T_{2}^{\prime} A_{4}
\end{array}\right) .
$$

Proposition 3 shows that if we want to construct functions $F^{\prime}$ which are CCZ-equivalent to a function $F$ and EA-inequivalent to both $F$ and $F^{-1}$ (if $F^{-1}$ exists), then we have to use a linear function $L_{1}(x, y)$ depending on both variables. However, this condition is not sufficient as the following example shows.

Example 1 Let $m=2 n+1$ and $s \equiv n[\bmod 2]$. Then the linear function

$$
\mathcal{L}(x, y)=\left(x+\operatorname{tr}(x)+\sum_{j=0}^{n-s} y^{2^{2 j+s}}, y+\operatorname{tr}(x)\right)
$$

is a permutation on $\mathbb{F}_{2^{m}}^{2}$ since the kernel of $\mathcal{L}$ is $\{(0,0)\}$ (this can be checked by considering the cases $\operatorname{tr}(x)=0$ and $\operatorname{tr}(x)=1$ ). For the Gold AB function $x^{3}$ the function

$$
F_{1}(x)=x+\operatorname{tr}(x)+\sum_{j=0}^{n-s}\left(x^{3}\right)^{2^{2 j+s}}
$$

is a permutation on $\mathbb{F}_{2^{m}}$. Indeed, denoting $L(y)=\sum_{j=0}^{n-s} y^{2^{2 j+s}}$ we have $L\left(y+y^{2}\right)=y+\operatorname{tr}(y)$ and $L\left((y+1)^{3}\right)=L\left(y^{3}\right)+y+\operatorname{tr}(y)+1$ since $L(1)=1$. Thus $F_{1}(x)=L\left((x+1)^{3}\right)+1$ and $F_{1}$ is a permutation if $L$ is bijective. The equation $z=L(y)$ implies $z+z^{2}=y+\operatorname{tr}(y)$\\
and $\operatorname{tr}(z)=\operatorname{tr}(y)$. Therefore, $L$ is a permutation and $L^{-1}(x)=x+x^{2}+\operatorname{tr}(x), F_{1}^{-1}(x)= \left[L^{-1}(x+1)\right]^{\frac{1}{3}}+1$. Finally, we get

$$
\begin{aligned}
F^{\prime}(x) & =F_{2} \circ F_{1}^{-1}(x)=\left(\left[L^{-1}(x+1)\right]^{\frac{1}{3}}+1\right)^{3}+\operatorname{tr}\left(\left[L^{-1}(x+1)\right]^{\frac{1}{3}}+1\right) \\
& =L^{-1}(x+1)+\left[L^{-1}(x+1)\right]^{\frac{2}{3}}+\left[L^{-1}(x+1)\right]^{\frac{1}{3}}+\operatorname{tr}\left(\left[L^{-1}(x+1)\right]^{\frac{1}{3}}\right) \\
& =L^{-1}(x+1)+L^{-1}\left(\left[L^{-1}(x+1)\right]^{\frac{1}{3}}\right)
\end{aligned}
$$

Thus, both functions $L_{1}$ and $L_{2}$ depend on two variables, but the function $F^{\prime}$ is EAequivalent to the inverse of $x^{3}$.\\
This example can be generalized for any Gold AB function by replacing $L^{-1}(x)=x+x^{2^{i}}+ \operatorname{tr}(x)$ and $x^{3}$ by $x^{2^{i}+1}$.

If we want to classify all functions CCZ-equivalent to a given one $F$, then we should characterize all permutations of the form $L \circ F+L^{\prime}$, where $L, L^{\prime}$ are linear. Indeed, we need to know which linear mapping $L_{1}: \mathbb{F}_{2}^{2 m} \rightarrow \mathbb{F}_{2}^{m}$ is such that the function $F_{1}(x)=L_{1}(x, F(x))$ is a permutation. Clearly, $L_{1}$ can be written uniquely in the form $L_{1}(x, y)=L(y)+L^{\prime}(x)$. If $F_{1}$ is a permutation then there exists a linear function $L_{2}(x, y)$ such that the linear function $\left(L_{1}, L_{2}\right)(x, y)$ is a permutation too. Indeed, $L_{1}(x, F(x))$ being a permutation, $L_{1}$ is onto and the kernel of $L_{1}$ has then dimension $m$. We can take for $L_{2}$ any linear permutation between $\operatorname{Ker}\left(L_{1}\right)$ and $\mathbb{F}_{2}^{m}$, extended to $\mathbb{F}_{2}^{2 m}$ by $L_{2}(x+y)=L_{2}(x)$ for all $x \in \operatorname{Ker}\left(L_{1}\right)$, $y \in E$, where $E$ is an $m$-dimensional subspace of $F_{2}^{2 m}$ such that $E \oplus \operatorname{Ker}\left(L_{1}\right)=\mathbb{F}_{2}^{2 m}$. Conversely, any linear function $L_{2}$ such that $\left(L_{1}, L_{2}\right)$ is a permutation has this form. Indeed, $L_{2}$ being onto, it has also an $m$-dimensional kernel, and ( $L_{1}, L_{2}$ ) being one to one, the kernels of $L_{1}$ and $L_{2}$ have trivial intersection. This proves that $\operatorname{Ker}\left(L_{1}\right) \oplus \operatorname{Ker}\left(L_{2}\right)=\mathbb{F}_{2}^{2 m}$ and that we can take $E=\operatorname{Ker}\left(L_{2}\right)$.

The following proposition gives a sufficient condition for a function to be EA-inequivalent to power functions.

Proposition 4 Let $F$ be a function from $\mathbb{F}_{2^{m}}$ to itself. If there exists an element $c \in \mathbb{F}_{2^{m}}^{*}$ such that $d^{\circ}(\operatorname{tr}(c F)) \notin\left\{0,1, d^{\circ}(F)\right\}$, then $F$ is EA-inequivalent to power functions.

Proof. It is proven in 8 that for any power function $x^{d}$ and for any $c \in \mathbb{F}_{2^{m}}$, either the function $\operatorname{tr}\left(c x^{d}\right)$ completely vanishes or it has exactly the algebraic degree $w_{2}(d)$. Thus, for any function $F$ which is affine equivalent to a power function, we have $d^{\circ}(\operatorname{tr}(c F)) \in \left\{0, d^{\circ}(F)\right\}, c \in \mathbb{F}_{2^{m}}^{*}$. Therefore, if $F$ is EA-equivalent to a power function then $d^{\circ}(\operatorname{tr}(c F)) \in \left\{0,1, d^{\circ}(F)\right\}$, for every $c \in \mathbb{F}_{2^{m}}^{*}$.

\section*{4 CCZ-equivalence and the Gold functions}
In the propositions below we give a characterization of permutations $L \circ F+L^{\prime}$ when $F$ is a Gold function. This characterization is not complete but it is useful for constructions of new APN and AB polynomials.

Proposition 5 Let $F: \mathbb{F}_{2^{m}} \rightarrow \mathbb{F}_{2^{m}}, F(x)=L\left(x^{2^{i}+1}\right)+L^{\prime}(x)$, where $L, L^{\prime}$ are linear and $\operatorname{gcd}(m, i)=1$. Then $F$ is a permutation if and only if, for every $u \neq 0$ in $\mathbb{F}_{2^{m}}$ and every $v$ such that $\operatorname{tr}(v)=\operatorname{tr}(1)$, the condition $L\left(u^{2^{i}+1} v\right) \neq L^{\prime}(u)$ holds.

Proof. For any $u \in \mathbb{F}_{2^{m}}^{*}$ we have

$$
\begin{aligned}
& F(x)+F(x+u)=L\left(x^{2^{i}+1}\right)+L^{\prime}(x)+L\left((x+u)^{2^{i}+1}\right)+L^{\prime}(x+u) \\
& =L\left(u x^{2^{i}}+u^{2^{i}} x+u^{2^{i}+1}\right)+L^{\prime}(u)=L\left(u^{2^{i}+1}\left((x / u)^{2^{i}}+x / u+1\right)\right)+L^{\prime}(u) .
\end{aligned}
$$

When $x$ ranges over $\mathbb{F}_{2^{m}}$ then $(x / u)^{2^{i}}+x / u+1$ ranges over $\left\{v \in \mathbb{F}_{2^{m}}: \operatorname{tr}(v)=\operatorname{tr}(1)\right\}$, since $\operatorname{gcd}(m, i)=1$. Hence, $F$ is injective (i.e. is a permutation) if and only if $L\left(u^{u^{i}+1} v\right) \neq L^{\prime}(u)$ for every $u \neq 0$ and every $v$ such that $\operatorname{tr}(v)=\operatorname{tr}(1)$.

Remark 3 If $m$ is even, then, up to affine equivalence and without loss of generality, we can consider only functions of the type $L\left(x^{2^{i}+1}\right)+x$. Indeed, if $F(x)=L\left(x^{2^{i}+1}\right)+L^{\prime}(x)$ is a permutation on $\mathbb{F}_{2^{m}}$ and $m$ is even, then it follows from Proposition 5 that $L^{\prime}$ must be a permutation (take $v=0$ ). Then the function $F^{\prime}(x)=L^{\prime-1} \circ F(x)=L^{\prime-1} \circ L\left(x^{2^{i}+1}\right)+x$ is a permutation if and only if $F$ is a permutation. Moreover if a function $L_{2}: \mathbb{F}_{2^{m}}^{2} \rightarrow \mathbb{F}_{2^{m}}$ is such that the function $\left(L(y)+L^{\prime}(x), L_{2}(x, y)\right)$ is a permutation on $\mathbb{F}_{2^{m}}^{2}$, then obviously $\left(L^{\prime-1} \circ L(y)+x, L_{2}(x, y)\right)$ is a permutation too; note also that $(L(y)+x, y)$ is a permutation for any linear function $L$. Thus the function $x^{2^{i}+1}$ is CCZ-equivalent to the functions $F_{2} \circ F^{\prime-1}$ and $F_{2} \circ F^{-1}$, where $F_{2}(x)=L_{2}\left(x, x^{2^{i}+1}\right)$. We have $F_{2} \circ F^{\prime-1}(x)=F_{2} \circ F^{-1} \circ L^{\prime}(x)$. Therefore, $F_{2} \circ F^{\prime-1}$ and $F_{2} \circ F^{-1}$ are affine equivalent.

Proposition 6 Let $F: \mathbb{F}_{2^{m}} \rightarrow \mathbb{F}_{2^{m}}, F(x)=L\left(x^{2^{i}+1}\right)+x$, where $L$ is linear, $m$ even and $\operatorname{gcd}(m, i)=1$. Let $L^{*}$ be the adjoint operator of $L$. Then $F$ is a permutation if and only if, for every $v \in \mathbb{F}_{2^{m}}$ such that $L^{*}(v) \neq 0$, there exists $u \in \mathbb{F}_{2^{m}}$ such that $L^{*}(v)=u^{2^{i}+1}$ and $\operatorname{tr}_{m / 2}\left(\frac{v}{u}\right) \neq 0$, where $\operatorname{tr}_{m / 2}$ is the trace function from $\mathbb{F}_{2^{m}}$ to $\mathbb{F}_{2^{2}}$.

Proof. The function $F$ is a permutation if and only if, for every $v \neq 0$, the Boolean function $\operatorname{tr}\left(v\left(L\left(x^{2^{i}+1}\right)+x\right)\right)$ is balanced (see e.g. 10). Let $L^{*}$ be the adjoint operator of $L$, then we have

$$
\operatorname{tr}\left(v\left(L\left(x^{2^{i}+1}\right)+x\right)\right)=\operatorname{tr}\left(L^{*}(v) x^{2^{i}+1}+v x\right) .
$$

If $L^{*}(v)=0$, then the function $\operatorname{tr}\left(v\left(L\left(x^{2^{i}+1}\right)+x\right)\right)$ is balanced. If $L^{*}(v) \neq 0$, the quadratic function $\operatorname{tr}\left(L^{*}(v) x^{2^{i}+1}+v x\right)$ has associated symplectic form:

$$
\varphi(x, y)=\operatorname{tr}\left(L^{*}(v) x^{2^{i}} y+L^{*}(v) x y^{2^{i}}\right)=\operatorname{tr}\left(\left(L^{*}(v) x^{2^{i}}+\left(L^{*}(v) x\right)^{2^{m-i}}\right) y\right),
$$

which has kernel :

$$
\begin{gathered}
\mathcal{E}=\left\{x \in \mathbb{F}_{2^{m}}: L^{*}(v) x^{2^{i}}+\left(L^{*}(v) x\right)^{2^{m-i}}=0\right\}= \\
=\left\{x \in \mathbb{F}_{2^{m}}: L^{*}(v)^{2^{i}} x^{2^{2 i}}+L^{*}(v) x=0\right\}=\{0\} \cup\left\{x \in \mathbb{F}_{2^{m}}: L^{*}(v)^{2^{i}-1} x^{2^{2 i}-1}=1\right\} .
\end{gathered}
$$

A quadratic function is balanced if and only if its restriction to the kernel of its associated symplectic form is not constant [8, 9]. The restriction of $\operatorname{tr}\left(L^{*}(v) x^{2^{i}+1}\right)$ to $\mathcal{E}$ is null. Indeed, $L^{*}(v)^{2^{i}-1} x^{2^{2 i}-1}=1$ implies that the order of $L^{*}(v) x^{2^{i}+1}$ divides $2^{i}-1$ and since $\operatorname{gcd}(i, m)=1$ then $L^{*}(v) x^{2^{i}+1} \in \mathbb{F}_{2}$ and the trace function is null on $\mathbb{F}_{2}$, since $m$ is even. Hence, the function $L\left(x^{2^{i}+1}\right)+x$ is a permutation if and only if every $v$ such that $L^{*}(v) \neq 0$ lies outside the dual of $\{0\} \cup\left\{x \in \mathbb{F}_{2^{m}}: L^{*}(v)^{2^{i}-1} x^{2^{2 i}-1}=1\right\}$. Equivalently, the function $L\left(x^{2^{i}+1}\right)+x$ is a permutation if and only if, for every $v$ such that $L^{*}(v) \neq 0$ the following two conditions hold:

\begin{enumerate}
  \item $L^{*}(v)^{2^{i}-1}$ belongs to $\left\{x^{2^{2 i}-1}: x \in \mathbb{F}_{2^{m}}\right\}$ (otherwise, $\mathcal{E}$ would be trivial), say $L^{*}(v)^{2^{i}-1}= u^{2^{2 i}-1}$, i.e. $L^{*}(v)=u^{2^{i}+1}$ (since the mapping $y \rightarrow y^{2^{i}-1}$ is a permutation); in this case $\mathcal{E}= \{0\} \cup\left\{x \in \mathbb{F}_{2^{m}}:(u x)^{2^{2 i}-1}=1\right\}=\frac{1}{u}\left\{y \in \mathbb{F}_{2^{m}}: y^{2^{2 i}}=y\right\}=\frac{1}{u} \mathbb{F}_{2^{j}}$, where $j=\operatorname{gcd}(2 i, m)=2$, hence $\mathcal{E}=\frac{1}{u} \mathbb{F}_{4}$;
  \item $v$ lies outside the dual of $\mathcal{E}$, that is, $L(v x) \neq 0$ for some $x \in \mathcal{E}$.
\end{enumerate}

For every $z \in \mathbb{F}_{2^{m}}$ and every $y \in \mathbb{F}_{2^{2}}$ we have

$$
\operatorname{tr}\left(z \frac{1}{u} y\right)=\operatorname{tr}_{2}\left(\operatorname{tr}_{m / 2}\left(\frac{z}{u} y\right)\right)=\operatorname{tr}_{2}\left(y \operatorname{tr}_{m / 2}\left(\frac{z}{u}\right)\right) .
$$

Hence, the function $L\left(x^{2^{i}+1}\right)+x$ is a permutation if and only if every $v$ such that $L^{*}(v) \neq 0$ satisfies $L^{*}(v)=u^{2^{i}+1}$ for some $u$ and $\operatorname{tr}_{m / 2}\left(\frac{v}{u}\right) \neq 0$.

\section*{5 New cases of AB and APN functions}
The next theorems show that CCZ-equivalent functions are not necessarily EA-equivalent (including the equivalence to the inverse). They lead to infinite classes of almost bent and almost perfect nonlinear polynomials, which are EA-inequivalent to power functions.

For the function $F^{\prime}: \mathbb{F}_{2^{m}} \rightarrow \mathbb{F}_{2^{m}}, F(x)=x^{2^{i}+1}, \operatorname{gcd}(m, i)=1$, and for any $a \in \mathbb{F}_{2^{m}}^{*}$ the set

$$
H_{a}=\left\{F(x+a)+F(x): x \in \mathbb{F}_{2^{m}}\right\}= \begin{cases}\left\{x \in \mathbb{F}_{2^{m}}: \operatorname{tr}\left(a^{-\left(2^{i}+1\right)} x\right)=1\right\} & \text { if } m \text { is odd } \\ \left\{x \in \mathbb{F}_{2^{m}}: \operatorname{tr}\left(a^{-\left(2^{i}+1\right)} x\right)=0\right\} & \text { if } m \text { is even }\end{cases}
$$

is a linear hyperplane when $m$ is even and a complement of a linear hyperplane when $m$ is odd (see 11 or Proposition 5).\\
We can use Proposition 1 . For any $a \in \mathbb{F}_{2^{m}}^{*}$ the set

$$
V^{\prime}= \begin{cases}\left(0, \mathbb{F}_{2^{m}} \backslash H_{a}\right) \cup\left(a, \mathbb{F}_{2^{m}} \backslash H_{a}\right) & \text { if } m \text { is odd } \\ \left(0, H_{a}\right) \cup\left(a, \mathbb{F}_{2^{m}} \backslash H_{a}\right) & \text { if } m \text { is even }\end{cases}
$$

is a subgroup of $\mathbb{F}_{2^{m}}^{2}$ such that $G_{F}$ is transversal to $V^{\prime}$ since $A_{F} \cap V^{\prime}=\emptyset$.\\
Theorem 1 The function $F^{\prime}: \mathbb{F}_{2^{m}} \rightarrow \mathbb{F}_{2^{m}}, m>3$ odd,

$$
F^{\prime}(x)=x^{2^{i}+1}+\left(x^{2^{i}}+x\right) \operatorname{tr}\left(x^{2^{i}+1}+x\right), \quad \operatorname{gcd}(m, i)=1,
$$

is an $A B$ function which is $E A$-inequivalent to any power function.

Proof. It is easy (but lengthy) to prove that for $a \in \mathbb{F}_{2^{m}}^{*}$ the linear mapping

$$
\mathcal{L}(x, y)=\left(x+a \operatorname{tr}\left(a^{-1} x\right)+a \operatorname{tr}\left(a^{-\left(2^{i}+1\right)} y\right), y+a^{2^{i}+1} \operatorname{tr}\left(a^{-\left(2^{i}+1\right)} y\right)+a^{2^{i}+1} \operatorname{tr}\left(a^{-1} x\right)\right)
$$

is an involution. We have

$$
\mathcal{L}(0, y)=\left\{\begin{array}{ll}
\left(a, y+a^{2^{i}+1}\right) & \text { if } y \in H_{a} \\
(0, y) & \text { if } y \in \mathbb{F}_{2^{m}} \backslash H_{a}
\end{array} .\right.
$$

Since $\operatorname{tr}\left(a^{-\left(2^{i}+1\right)} a^{2^{i}+1}\right)=1$, then $a^{2^{i}+1} \in H_{a}$. The set $H_{a}$ is a complement of a linear hyperplane, hence the sum of any two elements from $H_{a}$ belongs to $\mathbb{F}_{2^{m}} \backslash H_{a}$. Therefore, $\mathcal{L}^{-1}(V)=\mathcal{L}(V)=V^{\prime}$ and by Proposition 1 the function $F$ is CCZ-equivalent to $F_{2} \circ F_{1}^{-1}$, where

$$
\begin{aligned}
& F_{1}(x)=L_{1}(x, F(x))=x+a \operatorname{tr}(x / a)+a \operatorname{tr}\left((x / a)^{2^{i}+1}\right), \\
& F_{2}(x)=L_{2}(x, F(x))=x^{2^{i}+1}+a^{2^{i}+1} \operatorname{tr}\left((x / a)^{2^{i}+1}\right)+a^{2^{i}+1} \operatorname{tr}(x / a) .
\end{aligned}
$$

The function $F_{1}$ is an involution (this proves again that it is bijective):

$$
\begin{aligned}
F_{1} \circ F_{1}(x) & =x+a \operatorname{tr}(x / a)+a \operatorname{tr}\left((x / a)^{2^{i}+1}\right)+a \operatorname{tr}\left(a^{-1}(x+a \operatorname{tr}(x / a)\right. \\
& \left.\left.+a \operatorname{tr}\left((x / a)^{2^{i}+1}\right)\right)\right)+a \operatorname{tr}\left(a^{-\left(2^{i}+1\right)}\left(x+a \operatorname{tr}(x / a)+a \operatorname{tr}\left((x / a)^{2^{i}+1}\right)\right)^{2^{i}+1}\right) \\
& =x+3 a \operatorname{tr}(x / a)+2 a \operatorname{tr}\left((x / a)^{2^{i}+1}\right)+a \operatorname{tr}\left(a ^ { - ( 2 ^ { i } + 1 ) } \left(x^{2^{i}+1}+\left(a x^{2^{i}}+a^{2^{i}} x\right.\right.\right. \\
& \left.\left.\left.+a^{2^{i}+1}\right)\left(\operatorname{tr}(x / a)+\operatorname{tr}\left((x / a)^{2^{i}+1}\right)\right)\right)\right)=x+a \operatorname{tr}(x / a)+a \operatorname{tr}\left((x / a)^{2^{i}+1}\right) \\
& +a \operatorname{tr}\left((x / a)^{2^{i}}+(x / a)+1\right)\left(\operatorname{tr}(x / a)+\operatorname{tr}\left((x / a)^{2^{i}+1}\right)\right)=x,
\end{aligned}
$$

since

$$
\begin{aligned}
& \left(x+a \operatorname{tr}(x / a)+a \operatorname{tr}\left((x / a)^{2^{i}+1}\right)\right)^{2^{i}+1}=x^{2^{i}+1}+\left(a x^{2^{i}}+a^{2^{i}} x+a^{2^{i}+1}\right)\left(\operatorname{tr}(x / a)+\operatorname{tr}\left((x / a)^{2^{i}+1}\right)\right) \\
& \quad+2 a^{2^{i}+1} \operatorname{tr}(x / a) \operatorname{tr}\left((x / a)^{2^{i}+1}\right)=x^{2^{i}+1}+\left(a x^{2^{i}}+a^{2^{i}} x+a^{2^{i}+1}\right)\left(\operatorname{tr}(x / a)+\operatorname{tr}\left((x / a)^{2^{i}+1}\right)\right)
\end{aligned}
$$

Therefore,

$$
\begin{aligned}
F_{2} \circ F_{1}^{-1}(x) & =\left(x+a \operatorname{tr}(x / a)+a \operatorname{tr}\left((x / a)^{2^{i}+1}\right)\right)^{2^{i}+1}+a^{2^{i}+1} \operatorname{tr}\left(a^{-\left(2^{i}+1\right)}(x+a \operatorname{tr}(x / a)\right. \\
& \left.\left.+a \operatorname{tr}\left((x / a)^{2^{i}+1}\right)\right)^{2^{i}+1}\right)+a^{2^{i}+1} \operatorname{tr}\left(a^{-1}\left(x+a \operatorname{tr}(x / a)+a \operatorname{tr}\left((x / a)^{2^{i}+1}\right)\right)\right) \\
& =x^{2^{i}+1}+\left(a x^{2^{i}}+a^{2^{i}} x+a^{2^{i}+1}\right)\left(\operatorname{tr}(x / a)+\operatorname{tr}\left((x / a)^{2^{i}+1}\right)\right)+a^{2^{i}+1} \operatorname{tr}\left((x / a)^{2^{i}}\right. \\
& +(x / a)+1)\left(\operatorname{tr}(x / a)+\operatorname{tr}\left((x / a)^{2^{i}+1}\right)\right)+2 a^{2^{i}+1} \operatorname{tr}(x / a)+2 a^{2^{i}+1} \operatorname{tr}\left((x / a)^{2^{i}+1}\right) \\
& =a^{2^{i}+1}\left[(x / a)^{2^{i}+1}+\left((x / a)^{2^{i}}+(x / a)\right) \operatorname{tr}\left((x / a)+(x / a)^{2^{i}+1}\right)\right]=a^{2^{i}+1} F^{\prime}(x / a),
\end{aligned}
$$

where $F^{\prime}(x)=x^{2^{i}+1}+\left(x^{2^{i}}+x\right) \operatorname{tr}\left(x+x^{2^{i}+1}\right)$. By Proposition 2 the function $F^{\prime}$ is AB .\\
It is not difficult to see that the algebraic degree of $F^{\prime}$ is 3 for $m>3$. On the other hand $\operatorname{tr}\left(F^{\prime}(x)\right)=\operatorname{tr}\left(x^{2^{i}+1}\right)$ and $d^{\circ}\left(\operatorname{tr}\left(F^{\prime}(x)\right)\right)=2$. It follows from Proposition 4 that $F^{\prime}$ is EA-inequivalent to any power function.

Remark 4 It was conjectured in 11 that any AB function is EA-equivalent to a permutation. The AB function from Theorem 1 serves as a counterexample for this conjecture. It was checked by the help of a computer, that for no linear function $L$ on $\mathbb{F}_{2^{5}}$ the sum $F^{\prime}+L$ is a permutation for the AB function $F^{\prime}(x)=x^{2^{i}+1}+\left(x^{2^{i}}+x\right) \operatorname{tr}\left(x^{2^{i}+1}+x\right), \operatorname{gcd}(5, i)=1$. Thus, $F^{\prime}$ is EA-inequivalent to any permutation but it is CCZ-equivalent to the permutation $x^{2^{i}+1}$. For $m$ even the existence of APN permutations on $\mathbb{F}_{2^{m}}$ is an open problem. It is shown by Nyberg that there exist no quadratic APN permutations when $m$ is even [27. But as we have seen the nonexistence of permutations EA-equivalent to the Gold functions does not mean yet that there exist no permutations CCZ-equivalent to them.

Theorem 2 The function $F^{\prime}: \mathbb{F}_{2^{m}} \rightarrow \mathbb{F}_{2^{m}}, m \geq 4$ even,

$$
F^{\prime}(x)=x^{2^{i}+1}+\left(x^{2^{i}}+x+1\right) \operatorname{tr}\left(x^{2^{i}+1}\right), \quad \operatorname{gcd}(m, i)=1,
$$

is an APN function which is EA-inequivalent to any power function.\\
Proof. For $a \in \mathbb{F}_{2^{m}}^{*}$, the linear mapping

$$
\mathcal{L}(x, y)=\left(L_{1}, L_{2}\right)(x, y)=\left(x+a \operatorname{tr}\left(a^{-\left(2^{i}+1\right)} y\right), y\right)
$$

is an involution and, obviously, $\mathcal{L}(V)=V^{\prime}$. Thus, by Proposition 1 the function $F$ is CCZ-equivalent to $F_{2} \circ F_{1}^{-1}$, where

$$
\begin{aligned}
& F_{1}(x)=L_{1}(x, F(x))=x+a \operatorname{tr}\left((x / a)^{2^{i}+1}\right) \\
& F_{2}(x)=L_{2}(x, F(x))=x^{2^{i}+1}
\end{aligned}
$$

The function $F_{1}$ is an involution:

$$
\begin{aligned}
F_{1} \circ F_{1}(x) & =x+a \operatorname{tr}\left((x / a)^{2^{i}+1}\right)+a \operatorname{tr}\left(a ^ { - ( 2 ^ { i } + 1 ) } \left(x^{2^{i}+1}+a x^{2^{i}} \operatorname{tr}\left((x / a)^{2^{i}+1}\right)\right.\right. \\
& \left.\left.+a^{2^{i}} x \operatorname{tr}\left((x / a)^{2^{i}+1}\right)+a^{2^{i}+1} \operatorname{tr}\left((x / a)^{2^{i}+1}\right)\right)\right) \\
& =x+2 a \operatorname{tr}\left((x / a)^{2^{i}+1}\right)+a \operatorname{tr}\left((x / a)^{2^{i}}+(x / a)+1\right) \operatorname{tr}\left((x / a)^{2^{i}+1}\right)=x
\end{aligned}
$$

We have

$$
\begin{aligned}
F_{2} \circ F_{1}^{-1}(x) & =\left(x+a \operatorname{tr}\left((x / a)^{2^{i}+1}\right)\right)^{2^{i}+1}=x^{2^{i}+1}+a x^{2^{i}} \operatorname{tr}\left((x / a)^{2^{i}+1}\right) \\
& +a^{2^{i}} x \operatorname{tr}\left((x / a)^{2^{i}+1}\right)+a^{2^{i}+1} \operatorname{tr}\left((x / a)^{2^{i}+1}\right) \\
& =a^{2^{i}+1}\left[(x / a)^{2^{i}+1}+\left((x / a)^{2^{i}}+(x / a)+1\right) \operatorname{tr}\left((x / a)^{2^{i}+1}\right)\right]
\end{aligned}
$$

The function $F_{2} \circ F_{1}^{-1}$ is EA-equivalent to the function

$$
F^{\prime}(x)=x^{2^{i}+1}+\left(x^{2^{i}}+x+1\right) \operatorname{tr}\left(x^{2^{i}+1}\right) .
$$

Hence, $F^{\prime}$ is CCZ-equivalent to $F$ and it is APN by Proposition 2. The algebraic degree of $F^{\prime}$ is 3 and $d^{\circ}\left(\operatorname{tr}\left(F^{\prime}\right)\right)=2$ as $\operatorname{tr}\left(F^{\prime}(x)\right)=\operatorname{tr}\left(x^{2^{i}+1}\right)$. Therefore, $F^{\prime}$ is EA-inequivalent to

Note that the proofs of Theorems 1 and 2 do not depend on the condition $\operatorname{gcd}(i, m)=1$. When $\operatorname{gcd}(i, m)=s$ then the functions $F^{\prime}$ have the same differential and linear properties as $x^{2^{i}+1}$ and, therefore, if $m / s$ is odd they can be considered as the first polynomials with three valued Walsh spectrum $\left\{0, \pm 2^{\frac{m+s}{2}}\right\}$, which are EA-inequivalent to power functions.

Theorem 3 The function $F^{\prime}: \mathbb{F}_{2^{m}} \rightarrow \mathbb{F}_{2^{m}}, m$ divisible by 6,

$$
F^{\prime}(x)=\left[x+\operatorname{tr}_{m / 3}\left(x^{2\left(2^{i}+1\right)}+x^{4\left(2^{i}+1\right)}\right)+\operatorname{tr}(x) \operatorname{tr}_{m / 3}\left(x^{2^{i}+1}+x^{2^{2 i}\left(2^{i}+1\right)}\right)\right]^{2^{i}+1}
$$

with $\operatorname{gcd}(m, i)=1$, is an APN function.\\
Proof. The linear function $\mathcal{L}: \mathbb{F}_{2^{m}}^{2} \rightarrow \mathbb{F}_{2^{m}}^{2}$,

$$
\mathcal{L}(x, y)=\left(L_{1}, L_{2}\right)(x, y)=\left(x+t r_{m / 3}\left(y^{2}+y^{4}\right), y\right)
$$

is a permutation since it has the kernel $\{(0,0)\}$. For $m$ divisible by 6 , the function

$$
F_{1}(x)=L_{1}(x, F(x))=x+\operatorname{tr}_{m / 3}\left(x^{2\left(2^{i}+1\right)}+x^{4\left(2^{i}+1\right)}\right)
$$

is a permutation. Indeed, by Proposition 5 the function $F_{1}$ is a permutation if for every $v \in \mathbb{F}_{2^{m}}$ such that $\operatorname{tr}(v)=0$ and every $u \in \mathbb{F}_{2^{m}}^{*}$ the condition $\operatorname{tr}_{m / 3}\left(\left(u^{2^{i}+1} v\right)^{2}+\left(u^{2^{i}+1} v\right)^{4}\right) \neq u$ holds. Obviously, for any $u \notin \mathbb{F}_{2^{3}}$ the condition is satisfied and when $u \in \mathbb{F}_{2^{3}}^{*}$ the condition $\operatorname{tr}_{m / 3}\left(\left(u^{2^{i}+1} v\right)^{2}+\left(u^{2^{i}+1} v\right)^{4}\right) \neq u$ is equivalent to $\left(u^{2^{i}+1} \operatorname{tr}_{m / 3}(v)\right)^{2}+\left(u^{2^{i}+1} \operatorname{tr}_{m / 3}(v)\right)^{4} \neq u$. Therefore, $F_{1}$ is a permutation if, for every $u, w \in \mathbb{F}_{2^{3}}^{*}, \operatorname{tr}_{3}(w)=0$ the condition $\left(u^{2^{i}+1} w\right)^{2}+ \left(u^{2^{i}+1} w\right)^{4} \neq u$ is satisfied and that was easily checked by a computer.

We show below that $F_{1}^{-1}=F_{1} \circ F_{1} \circ F_{1} \circ F_{1} \circ F_{1}$.\\
We denote

$$
T(x)=t r_{m / 3}\left(x^{2^{i}+1}\right)
$$

then

$$
F_{1}(x)=x+T(x)^{2}+T(x)^{4}
$$

Since every element of $\mathbb{F}_{8}$ is equal to its 8 power and the function $\operatorname{tr}_{m / 3}(x)$ is 0 on $\mathbb{F}_{8}$, then

$$
\begin{aligned}
& T \circ F_{1}(x)=\operatorname{tr}_{m / 3}\left[\left(x+\operatorname{tr}_{m / 3}\left(x^{2\left(2^{i}+1\right)}+x^{4\left(2^{i}+1\right)}\right)\right)^{2^{i}+1}\right] \\
= & \operatorname{tr}_{m / 3}\left[( x + \operatorname { t r } _ { m / 3 } ( x ^ { 2 ( 2 ^ { i } + 1 ) } + x ^ { 4 ( 2 ^ { i } + 1 ) } ) ) \left(x^{2^{i}}+\operatorname{tr}_{m / 3}\left(x^{2^{i+1}\left(2^{i}+1\right)}\right.\right.\right. \\
+ & \left.\left.\left.x^{2^{i+2}\left(2^{i}+1\right)}\right)\right)\right]=\operatorname{tr}_{m / 3}\left[x^{2^{i}+1}+x \operatorname{tr}_{m / 3}\left(x^{2^{i+1}\left(2^{i}+1\right)}+x^{2^{i+2}\left(2^{i}+1\right)}\right)\right. \\
+ & \left.x^{2^{i}} \operatorname{tr}_{m / 3}\left(x^{2\left(2^{i}+1\right)}+x^{4\left(2^{i}+1\right)}\right)\right]=\operatorname{tr}_{m / 3}(x) \operatorname{tr}_{m / 3}\left(x^{2^{s+1}\left(2^{i}+1\right)}+x^{s^{s+2}\left(2^{i}+1\right)}\right) \\
+ & \operatorname{tr}_{m / 3}\left(x^{2^{s}}\right) \operatorname{tr}_{m / 3}\left(x^{2\left(2^{i}+1\right)}+x^{4\left(2^{i}+1\right)}\right)+\operatorname{tr}_{m / 3}\left(x^{2^{i}+1}\right)
\end{aligned}
$$

where $s \equiv i \bmod 3$.\\
Therefore,

$$
\begin{aligned}
& \left.\left.F_{1} \circ F_{1}(x)=F_{1}(x)+\left[T \circ F_{1}(x)\right)\right]^{2}+\left[T \circ F_{1}(x)\right)\right]^{4} \\
= & x+2 \operatorname{tr}_{m / 3}\left(x^{2\left(2^{i}+1\right)}+x^{4\left(2^{i}+1\right)}\right)+\operatorname{tr}_{m / 3}\left(x^{2}\right) \operatorname{tr}_{m / 3}\left(x^{2^{s+2}\left(2^{i}+1\right)}\right. \\
+ & \left.x^{2^{s}\left(2^{i}+1\right)}\right)+\operatorname{tr}_{m / 3}\left(x^{2^{s+1}}\right) \operatorname{tr}_{m / 3}\left(x^{4\left(2^{i}+1\right)}+x^{2^{i}+1}\right) \\
+ & \operatorname{tr}_{m / 3}\left(x^{4}\right) \operatorname{tr}_{m / 3}\left(x^{2^{s}\left(2^{i}+1\right)}+x^{2^{s+1}\left(2^{i}+1\right)}\right)+\operatorname{tr}_{m / 3}\left(x^{2^{s+2}}\right) \operatorname{tr}_{m / 3}\left(x^{2^{i}+1}+x^{2\left(2^{i}+1\right)}\right)
\end{aligned}
$$

Considering separately the cases $s=1$ and $s=2$ we get

$$
\begin{aligned}
F_{1} \circ F_{1}(x) & =x+\operatorname{tr}_{m / 3}\left(x+x^{2}+x^{4}\right) \operatorname{tr}_{m / 3}\left(x^{2^{i}+1}+x^{2^{s}\left(2^{i}+1\right)}\right) \\
& =x+\operatorname{tr}(x) \operatorname{tr}_{m / 3}\left(x^{2^{i}+1}+x^{2^{s}\left(2^{i}+1\right)}\right)=x+\left(T(x)+T(x)^{2^{s}}\right) \operatorname{tr}(x)
\end{aligned}
$$

Like this we get

$$
\begin{gathered}
F_{1} \circ F_{1} \circ F_{1} \circ F_{1} \circ F_{1}(x)=x+T(x)^{2}+T(x)^{4}+\operatorname{tr}(x)\left(T(x)+T(x)^{2^{2 s}}\right) \\
F_{1} \circ F_{1} \circ F_{1} \circ F_{1} \circ F_{1} \circ F_{1}(x)=x
\end{gathered}
$$

Thus,

$$
\begin{aligned}
F^{\prime}(x) & =F_{2} \circ F_{1}^{-1}(x)=\left[x+T(x)^{2}+T(x)^{4}+\operatorname{tr}(x)\left(T(x)+T(x)^{2^{2 s}}\right)\right]^{2^{i}+1}=x^{2^{i}+1}+T(x)^{2^{s}+1} \\
& +\operatorname{tr}\left(x^{2^{i}+1}\right) T(x)^{2^{2 s}}+\operatorname{tr}(x)\left(T(x)+T(x)^{4}\right)+x \operatorname{tr}(x)\left(T(x)+T(x)^{2^{s}}\right) \\
& +x^{2^{i}} \operatorname{tr}(x)\left(T(x)+T(x)^{2^{2 s}}\right)+x\left(T(x)+T(x)^{2^{2 s}}\right)+x^{2^{i}}\left(T(x)^{2}+T(x)^{4}\right),
\end{aligned}
$$

where $F_{2}(x)=L_{2}(x, F(x))=x^{2^{i}+1}$.\\
The function $F^{\prime}$ has the algebraic degree 4, this can be shown by lengthy but uncomplicated calculations. Hence, $F^{\prime}$ is EA-inequivalent to other known APN functions since for $m$ divisible by 6 we have no known APN functions of algebraic degree 4 .

Let $F: \mathbb{F}_{2^{m}} \rightarrow \mathbb{F}_{2^{m}}, F(x)=x^{2^{i}+1}$ with $\operatorname{gcd}(m, i)=1, m$ odd and $n$ a divisor of $m$. Then $A_{F} \cap V^{\prime}=\emptyset$ for the subgroup

$$
V^{\prime}=\left\{(b, x): b \in \mathbb{F}_{2^{n}}, x \in \mathbb{F}_{2^{m}}, \operatorname{tr}_{m / n}(x)=0\right\}
$$

since if $(b, x) \in V^{\prime}$ then $\operatorname{tr}\left(b^{-\left(2^{i}+1\right)} x\right)=\operatorname{tr}_{n}\left(b^{-\left(2^{i}+1\right)} \operatorname{tr}_{m / n}(x)\right)=0$ and $(b, x) \notin A_{F}$. Hence, $G_{F}$ is transversal to $V^{\prime}$. Using the subgroup $V^{\prime}$ we construct an AB function $F^{\prime}$ given in the following theorem.\\
Theorem 4 The function $F^{\prime}: \mathbb{F}_{2^{m}} \rightarrow \mathbb{F}_{2^{m}}$, where $m$ is odd and divisible by $n, m \neq n$ and $\operatorname{gcd}(m, i)=1$,

$$
\begin{aligned}
F^{\prime}(x) & =x^{2^{i}+1}+\operatorname{tr}_{m / n}\left(x^{2^{i}+1}\right)+x^{2^{i}} \operatorname{tr}_{m / n}(x)+x \operatorname{tr}_{m / n}(x)^{2^{i}} \\
& +\left[\operatorname{tr}_{m / n}(x)^{2^{i}+1}+\operatorname{tr}_{m / n}\left(x^{2^{i}+1}\right)+\operatorname{tr}_{m / n}(x)\right]^{\frac{1}{2^{i}+1}}\left(x^{2^{i}}+\operatorname{tr}_{m / n}(x)^{2^{i}}+1\right) \\
& +\left[\operatorname{tr}_{m / n}(x)^{2^{i}+1}+\operatorname{tr}_{m / n}\left(x^{2^{i}+1}\right)+\operatorname{tr}_{m / n}(x)\right]^{2^{i}+1}\left(x+\operatorname{tr}_{m / n}(x)\right)
\end{aligned}
$$

is an $A B$ function which is $E A$-inequivalent to any power function.

Proof. Let $m$ be odd and divisible by $n$. Obviously, the linear function

$$
\mathcal{L}(x, y)=\left(L_{1}, L_{2}\right)(x, y)=\left(x+t r_{m / n}(x)+t r_{m / n}(y), y+t r_{m / n}(x)\right)
$$

is a permutation on $\mathbb{F}_{2^{m}}^{2}$ and

$$
\mathcal{L}^{-1}(x, y)=\left(x+t r_{m / n}(y), y+t r_{m / n}(x)+t r_{m / n}(y)\right)
$$

We have

$$
\begin{gathered}
F_{1}(x)=x+t r_{m / n}(x)+t r_{m / n}\left(x^{2^{i}+1}\right) \\
F_{2}(x)=x^{2^{i}+1}+t r_{m / n}(x)
\end{gathered}
$$

By Proposition 1, the function $F_{1}$ is a permutation since $\mathcal{L}^{-1}(V)=V^{\prime}$ and $G_{F}$ is transversal to $V^{\prime}$. We need the inverse of the function $F_{1}$ to construct $F^{\prime}=F_{2} \circ F_{1}^{-1}$.\\
For any fixed element $x \in \mathbb{F}_{2^{m}}$ we have

$$
y=x+t r_{m / n}(x)+t r_{m / n}\left(x^{2^{i}+1}\right)=x+u
$$

for some $u \in \mathbb{F}_{2^{n}}$, and, therefore, $x=y+u$. Then

$$
y=(y+u)+t r_{m / n}(y+u)+t r_{m / n}\left((y+u)^{2^{i}+1}\right)
$$

which yields


\begin{equation*}
u^{2^{i}+1}+u^{2^{i}} t r_{m / n}(y)+u\left(t r_{m / n}(y)\right)^{2^{i}}+t r_{m / n}\left(y^{2^{i}+1}\right)+t r_{m / n}(y)=0 . \tag{4}
\end{equation*}


If $\operatorname{tr}_{m / n}(y) \neq 0$ then we denote $v=u / t r_{m / n}(y)$ and we get

$$
v^{2^{i}+1}+v^{2^{i}}+v+\frac{t r_{m / n}\left(y^{2^{i}+1}\right)+t r_{m / n}(y)}{\left(t r_{m / n}(y)\right)^{2^{i}+1}}=0
$$

Since $v^{2^{i}+1}+v^{2^{i}}+v=(v+1)^{2^{i}+1}+1$ then

$$
v+1=\left[\frac{t r_{m / n}\left(y^{2^{i}+1}\right)+t r_{m / n}(y)}{\left(t r_{m / n}(y)\right)^{2^{i}+1}}+1\right]^{\frac{1}{2^{i}+1}}
$$

Replacing $v$ by $u / t r_{m / n}(y)$ we have

$$
u=\left[\left(t r_{m / n}(y)\right)^{2^{i}+1}+t r_{m / n}\left(y^{2^{i}+1}\right)+t r_{m / n}(y)\right]^{\frac{1}{2^{i}+1}}+t r_{m / n}(y)
$$

If $\operatorname{tr}_{m / n}(y)=0$ then from the equality (4) we get $u=\left[\operatorname{tr}_{m / n}\left(y^{2^{i}+1}\right)\right]^{\frac{1}{2^{i}+1}}$ and we observe that $u$ equals again $\left[\left(t r_{m / n}(y)\right)^{2^{i}+1}+t r_{m / n}\left(y^{2^{i}+1}\right)+t r_{m / n}(y)\right]^{\frac{1}{2^{i}+1}}+t r_{m / n}(y)$. Thus, in all cases, we have

$$
F_{1}^{-1}(y)=y+u=y+\left[\left(t r_{m / n}(y)\right)^{2^{i}+1}+t r_{m / n}\left(y^{2^{i}+1}\right)+t r_{m / n}(y)\right]^{\frac{1}{2^{i}+1}}+t r_{m / n}(y)
$$

and

$$
\begin{aligned}
F^{\prime}(x) & =F_{2} \circ F_{1}^{-1}(x)=\left[x+\left[\left(\operatorname{tr}_{m / n}(x)\right)^{2^{i}+1}+\operatorname{tr}_{m / n}\left(x^{2^{i}+1}\right)+\operatorname{tr}_{m / n}(x)\right]^{\frac{1}{2^{i}+1}}+\operatorname{tr}_{m / n}(x)\right]^{2^{i}+1} \\
& +\operatorname{tr}_{m / n}\left[x+\left[\left(\operatorname{tr}_{m / n}(x)\right)^{2^{i}+1}+\operatorname{tr}_{m / n}\left(x^{2^{i}+1}\right)+\operatorname{tr}_{m / n}(x)\right]^{\frac{1}{2^{i}+1}}+\operatorname{tr}_{m / n}(x)\right] \\
& =x^{2^{i}+1}+\operatorname{tr}_{m / n}\left(x^{2^{i}+1}\right)+\operatorname{tr}_{m / n}(x)+x^{2^{i}} \operatorname{tr}_{m / n}(x)+x\left(\operatorname{tr}_{m / n}(x)\right)^{2^{i}} \\
& +\left[\left(\operatorname{tr}_{m / n}(x)\right)^{2^{i}+1}+\operatorname{tr}_{m / n}\left(x^{2^{i}+1}\right)+\operatorname{tr}_{m / n}(x)\right]^{\frac{1}{2^{i}+1}}\left(x^{2^{i}}+\left(\operatorname{tr}_{m / n}(x)\right)^{2^{i}}+1\right) \\
& +\left[\left(\operatorname{tr}_{m / n}(x)\right)^{2^{i}+1}+\operatorname{tr}_{m / n}\left(x^{2^{i}+1}\right)+\operatorname{tr}_{m / n}(x)\right]^{\frac{2^{i}}{2^{i}+1}}\left(x+\operatorname{tr}_{m / n}(x)\right)
\end{aligned}
$$

We show below that the function $F^{\prime}$ has the algebraic degree $n+2$. It means that the number of functions CCZ-equivalent to a Gold AB function and EA-inequivalent to it is not smaller than the number of divisors of $m$.

The inverse of $x^{2^{i}+1}$ on $F_{2^{n}}$ is $x^{d}$, where

$$
d=\sum_{k=0}^{\frac{n-1}{2}} 2^{2 i k}
$$

and $x^{d}$ has the algebraic degree $\frac{n+1}{2}$ [26]. Obviously, $\left(\left(t_{m / n}(x)\right)^{2^{i}+1}+t r_{m / n}\left(x^{2^{i}+1}\right)+\right. \left.t r_{m / n}(x)\right)^{d}$ has the algebraic degree $n+1$ if and only if $\left(\left(t r_{m / n}(x)\right)^{2^{i}+1}+t r_{m / n}\left(x^{2^{i}+1}\right)\right)^{d}$ has this algebraic degree.

We assume that $m \neq n$ and $n \neq 1$ since when $m=n$ we get $F^{\prime}(x)=x^{\frac{1}{2^{i}+1}}+x$ and Theorem 1 gives the case $n=1$.

We have

$$
\operatorname{tr}_{m / n}(x)=\sum_{k=0}^{\frac{m}{n}-1} x^{2^{k n}}
$$

and

$$
\begin{gathered}
\left(\operatorname{tr}_{m / n}(x)\right)^{2^{i}+1}+t r_{m / n}\left(x^{2^{i}+1}\right)=\sum_{k=0}^{\frac{m}{n}-1} x^{2^{k n}} \sum_{k=0}^{\frac{m}{n}-1} x^{2^{k n+i}}+\sum_{k=0}^{\frac{m}{n}-1}\left(x^{2^{i}+1}\right)^{2^{k n}} \\
=\sum_{k, j=0}^{\frac{m}{n}-1} x^{2^{k n}+2^{j n+i}}+\sum_{k=0}^{\frac{m}{n}-1} x^{2^{k n}+2^{k n+i}}=\sum_{\substack{k, j=0 \\
k \neq j}}^{\frac{m}{n}-1} x^{k^{n}+2^{j n+i}}
\end{gathered}
$$

Note that we have

$$
\begin{gathered}
{\left[\left(\operatorname{tr}_{m / n}(x)\right)^{2^{i}+1}+t r_{m / n}\left(x^{2^{i}+1}\right)\right]^{2^{2 i}+1}=\sum_{\substack{k, j=0 \\
k \neq j}}^{\frac{m}{n}-1} x^{k^{n}+2^{j n+i}} \sum_{\substack{k, j=0 \\
k \neq j}}^{\frac{m}{n}-1} x^{2^{k n+2 i}+2^{j n+3 i}}} \\
=\sum_{\substack{k, j, s, t=0 \\
k \neq j, s \neq t}}^{\frac{m}{n}-1} x^{2^{k n}+2^{j n+i}+2^{s n+2 i}+2^{t n+3 i}}
\end{gathered}
$$

Similarly, we have


\begin{equation*}
\left(\left(t r_{m / n}(x)\right)^{2^{i}+1}+t r_{m / n}\left(x^{2^{i}+1}\right)\right)^{d}=\sum_{\left(k_{0}, \ldots, k_{n}\right) \in I} x^{\sum_{s=0}^{n} 2^{k_{s} n+s i}} \tag{5}
\end{equation*}


where $I=\left\{\left(l_{0}, \ldots, l_{n}\right): \quad 0 \leq l_{t} \leq \frac{m}{n}-1, \quad l_{2 t} \neq l_{2 t+1}\right\}$.\\
The equality $k_{s} n+s i=k_{t} n+t i$ is possible for $0 \leq s<t \leq n$ only for $s=0$ and $t=n$. Indeed, if $k_{s} n+s i=k_{t} n+t i$ then $\left(k_{s}-k_{t}\right) n=(t-s) i$. Since $\operatorname{gcd}(n, i)=1$ and $0 \leq t, s \leq n$ then $t=n, s=0$ and $k_{s}-k_{t}=i$.

For simplicity we consider now the equality (5) in the case $i=1$. In the sum $\sum_{s=0}^{n} 2^{k_{s} n+s}$ the largest possible item is $2^{\left(\frac{m}{n}-1\right) n+n}=2^{m}$. Therefore, when $k_{n} \neq \frac{m}{n}-1$ the sum is smaller than $2^{m}-1$. Besides, all items in the sum are different modulo $2^{m}-1$ except the case when $k_{0}=0$ and $k_{n}=\frac{m}{n}-1$ and in the cases where $k_{0}=k_{n}+1$. Therefore, when $k_{0}=k_{n}=1$ the number $\sum_{s=0}^{n} 2^{k_{s} n+s}$ has the weight $n+1$. On the other hand, when $k_{0}=k_{n}=1$ and $k_{1}=k_{n-1}=0$ the term

$$
x^{\sum_{s=0}^{n} 2^{k_{s} n+s}}
$$

does not vanish in (5). Indeed, if

$$
\sum_{s=0}^{n} 2^{k_{s} n+s} \equiv \sum_{p=0}^{n} 2^{t_{p} n+p} \quad \bmod \left(2^{m}-1\right)
$$

then we have only two possibilities:

\begin{enumerate}
  \item for any $s$ there exists $p$ such that $k_{s} n+s=t_{p} n+p$ (and vice versa). Then $\left(k_{s}-t_{p}\right) n=p-s$ and since $0 \leq s, p \leq n$ then $k_{0}=t_{n}+1, t_{0}=k_{n}+1$ and $k_{s}=t_{s}$ for $s \neq 0, n$. If $k_{0}=k_{n}=1$ then $t_{n}=0, t_{0}=2$. But in our case $k_{0}=k_{n}=1, t_{1}=k_{1}=0, t_{n-1}=k_{n-1}=0$ and, therefore, $t_{n} \neq 0$ since $t_{n-1} \neq t_{n}$.
  \item if $t_{n}=\frac{m}{n}-1$ then $k_{0}$ must be equal to 0 or $k_{n}=\frac{m}{n}-1$, but $k_{0}=k_{n}=1$.
\end{enumerate}

Thus, when $k_{0}=k_{n}=1$ and $k_{1}=k_{n-1}=0$ (for permissible $k_{s}, 1<s<n-1$ ) the term

$$
x^{\sum_{s=0}^{n} 2^{k_{s} n+s}}
$$

has the algebraic degree $n+1$ and it does not vanish in (5).\\
If $n \geq 5$ we can also take $k_{2}=1, k_{3}=k_{4}=0$ and then we get


\begin{equation*}
\sum_{s=0}^{n} 2^{k_{s} n+s}=2^{n}+2+2^{n+2}+2^{3}+2^{4}+\ldots+2^{2 n} \tag{6}
\end{equation*}


We have

$$
\begin{aligned}
& \left(\left(t r_{m / n}(x)\right)^{3}+t r_{m / n}\left(x^{3}\right)\right)^{d}\left(x^{2}+t r_{m / n}(x)^{2}\right)+\left(\left(t r_{m / n}(x)\right)^{3}+t r_{m / n}\left(x^{3}\right)\right)^{2 d}\left(x+t r_{m / n}(x)\right) \\
& \quad=\sum_{\left(k_{0}, \ldots, k_{n}\right) \in I} x^{\sum_{s=0}^{n} 2^{k_{s} n+s}} \sum_{1 \leq k \leq m / n-1} x^{2^{n k+1}}+\sum_{\left(k_{0}, \ldots, k_{n}\right) \in I} x^{\sum_{s=0}^{n} 2^{k_{s} n+s+1}} \sum_{1 \leq j \leq m / n-1} x^{2^{n j}}
\end{aligned}
$$


\begin{equation*}
=\sum_{\substack{\left(k_{0}, \ldots, k_{n}\right) \in I \\ 1 \leq k \leq m / n-1}} x^{2^{n k+1}+\sum_{s=0}^{n} 2^{k_{s} n+s}}+\sum_{\substack{\left(k_{0}, \ldots, k_{n}\right) \in I \\ 1 \leq j \leq m / n-1}} x^{2^{n j}+\sum_{s=0}^{n} 2^{k_{s} n+s+1}} . \tag{7}
\end{equation*}


We consider the item with the exponent


\begin{equation*}
2^{n}+2+2^{n+2}+2^{3}+2^{4}+\ldots+2^{2 n}+2^{n k+1} \tag{8}
\end{equation*}


from the first sum in (7). It is easy to see that $2^{n}+2+2^{n+2}+2^{3}+2^{4}+\ldots+2^{2 n}+2^{n k+1}<2^{m}$ since $k \leq m / n-1$. In this sum $n k+1=k_{s} n+s$ only if $s=1$. But then $k=k_{1}=0$ which is in contradiction with $1 \leq k$. Thus, the number given by this sum has the weight $n+2$. The item with the exponent (8) does not vanish. Indeed, if there is another item in the first sum of (7) with this exponent then

$$
2^{n}+2+2^{n+2}+2^{3}+2^{4}+\ldots+2^{2 n}+2^{n k+1}=2^{n j+1}+\sum_{s=0}^{n} 2^{k_{s} n+s}
$$

If $k=j$ then (6) is equal to another sum $\sum_{s=0}^{n} 2^{k_{s} n+s}$ and we already showed that it is impossible. If $k \neq j$ then $k_{1}=k$ and $j=0$ while $1 \leq j$.\\
Assume there exists an item in the second sum of (7) with the exponent (8) then

$$
2^{n}+2+2^{n+2}+2^{3}+2^{4}+\ldots+2^{2 n}+2^{n k+1}=2^{n j}+\sum_{s=0}^{n} 2^{k_{s} n+s+1}
$$

for some $j, 1 \leq j \leq m / n-1$ and $\left(k_{0}, \ldots, k_{n}\right) \in I$. We have $3=k_{s} n+s+1 \bmod m$ for some $s, 0 \leq s \leq n$. Then $k_{s} n=2-s$ or $k_{s} n=m-(s-2)$ and this is possible only if $k_{2}=0$ or $k_{2}=m / n$, but since $0 \leq k_{s} \leq m / n-1$, then $3=k_{s} n+s+1 \bmod m$ only if $k_{2}=0$. The same arguments show that $4=k_{s} n+s+1 \bmod m$ only if $k_{3}=0$ and that is in contradiction with the condition $k_{2 t} \neq k_{2 t+1}$. Therefore the item with the exponent (8) does not vanish in (7) and then it does not vanish in the sum presenting the function $F^{\prime}$. This completes the proof that $F^{\prime}$ has the algebraic degree $n+2$.

The algebraic degree of the function $\operatorname{tr}\left(F^{\prime}(x)\right)$ is not greater than $n+1$ since

$$
\begin{aligned}
\operatorname{tr}\left(F^{\prime}(x)\right) & =\operatorname{tr}\left(x^{2^{i}+1}+\operatorname{tr}_{m / n}\left(x^{2^{i}+1}\right)+\operatorname{tr}_{m / n}(x)+x^{2^{i}} \operatorname{tr}_{m / n}(x)+x\left(\operatorname{tr}_{m / n}(x)\right)^{2^{i}}\right) \\
& +\operatorname{tr}_{n}\left(\left[\left(\operatorname{tr}_{m / n}(x)\right)^{2^{i}+1}+\operatorname{tr}_{m / n}\left(x^{2^{i}+1}\right)+\operatorname{tr}_{m / n}(x)\right]^{\frac{1}{2^{i}+1}} \operatorname{tr}_{m / n}\left(x^{2^{i}}+\left(\operatorname{tr}_{m / n}(x)\right)^{2^{i}}+1\right)\right) \\
& +\operatorname{tr}_{n}\left(\left[\left(\operatorname{tr}_{m / n}(x)\right)^{2^{i}+1}+\operatorname{tr}_{m / n}\left(x^{2^{i}+1}\right)+\operatorname{tr}_{m / n}(x)\right]^{2^{i}+1} \operatorname{tr}_{m / n}\left(x+\operatorname{tr}_{m / n}(x)\right)\right) \\
& =\operatorname{tr}(x)+\operatorname{tr}\left(x^{2^{i}} \operatorname{tr}_{m / n}(x)\right)+\operatorname{tr}\left(x\left(\operatorname{tr}_{m / n}(x)\right)^{2^{i}}\right) \\
& +\operatorname{tr}\left(\left[\left(\operatorname{tr}_{m / n}(x)\right)^{2^{i}+1}+\operatorname{tr}_{m / n}\left(x^{2^{i}+1}\right)+\operatorname{tr}_{m / n}(x)\right]^{\frac{1}{2^{i}+1}}\right)
\end{aligned}
$$

On the other hand $d^{\circ}\left(\operatorname{tr}\left(F^{\prime}(x)\right)\right)$ is not 0 or 1 . Indeed, for any $x \in \mathbb{F}_{2^{n}}$ we have

$$
\operatorname{tr}\left(F^{\prime}(x)\right)=\operatorname{tr}(x)+2 \operatorname{tr}\left(x^{2^{i}+1}\right)+\operatorname{tr}\left(\left[2 x^{2^{i}+1}+x\right]^{\frac{1}{2^{i}+1}}\right)=\operatorname{tr}(x)+\operatorname{tr}\left(x^{\frac{1}{2^{i}+1}}\right)
$$

The function $\operatorname{tr}\left(x^{\frac{1}{2^{i}+1}}\right)$ has the algebraic degree $(n+1) / 2$. Indeed, $d^{\circ}\left(\operatorname{tr}\left(x^{\frac{1}{2^{i}+1}}\right)\right) \in\left\{0, d^{\circ}\left(x^{\frac{1}{2^{i}+1}}\right)\right\}$ and since $x^{\frac{1}{2^{i}+1}}$ is a permutation then $\operatorname{tr}\left(x^{\frac{1}{2^{i}+1}}\right)$ is not a constant and $d^{\circ}\left(\operatorname{tr}\left(x^{\frac{1}{2^{i}+1}}\right)\right)= d^{\circ}\left(x^{\frac{1}{2^{i}+1}}\right)=(n+1) / 2$. Hence, $\operatorname{tr}\left(F^{\prime}(x)\right.$ is not linear on $\mathbb{F}_{2^{n}}$ and then it cannot be linear on $\mathbb{F}_{2^{m}}$.

Thus, the function $F^{\prime}$ is EA-inequivalent to any power function by Proposition 4.

\section*{6 Conclusion}
We have shown that there exist APN and AB functions which are EA -inequivalent to power functions, and therefore, which are new, up to EA-equivalence. But these functions are CCZ-equivalent to the Gold functions. We leave two open problems:

\begin{itemize}
  \item finding classes of functions CCZ-equivalent to other known APN or AB functions, which would be EA-inequivalent to all known APN and AB functions (or even, inequivalent to power functions);
  \item finding classes of APN or AB functions which would be CCZ-inequivalent to all known APN and AB functions (or even, CCZ-inequivalent to power functions).
\end{itemize}

\section*{References}
[1] T. Bending, D. Fon-Der-Flaass. Crooked functions, bent functions and distanceregular graphs. Electron. J. Comb., 5(R34), 14, 1998.\\[0pt]
[2] T. Beth and C. Ding. On almost perfect nonlinear permutations. Advances in Cryptology-EUROCRYPT'93, Lecture Notes in Computer Science, 765, SpringerVerlag, New York, pp. 65-76, 1993.\\[0pt]
[3] E. Biham and A. Shamir. Differential Cryptanalysis of DES-like Cryptosystems. Journal of Cryptology, vol. 4, No.1, pp. 3-72, 1991.\\[0pt]
[4] L. Budaghyan, C. Carlet, A. Pott. New Constructions of Almost Perfect Nonlinear and Almost Bent Functions. Proceedings of the Workshop on Coding and Cryptography 2005, P. Charpin and $\varnothing$. Ytrehus eds, pp. 306-315, 2005.\\[0pt]
[5] A. Canteaut, P. Charpin and H. Dobbertin. A new characterization of almost bent functions. Fast Software Encryption 99, Lecture Notes in Computer Science 1636, L. Knudsen edt, pp. 186-200. Springer-Verlag, 1999.\\[0pt]
[6] A. Canteaut, P. Charpin and H. Dobbertin. Binary $m$-sequences with three-valued crosscorrelation: A proof of Welch's conjecture. IEEE Trans. Inform. Theory, 46 (1), pp. 4-8, 2000.\\[0pt]
[7] A. Canteaut, P. Charpin, H. Dobbertin. Weight divisibility of cyclic codes, highly nonlinear functions on $\mathbb{F}_{2^{m}}$, and crosscorrelation of maximum-length sequences. SIAM Journal on Discrete Mathematics, 13(1), pp. 105-138, 2000.\\[0pt]
[8] C. Carlet. Codes de Reed-Muller, codes de Kerdock et de Preparata, PhD thesis, Publication of LITP, Institut Blaise Pascal, Université Paris 6, 90.59, 1990.\\[0pt]
[9] C. Carlet. Boolean Functions for Cryptography and Error Correcting Codes. Chapter of the monography Boolean Methods and Models, Y. Crama and P. Hammer eds, Cambridge University Press, to appear (winter 2005-2006).\\[0pt]
[10] C. Carlet. Vectorial Boolean Functions for Cryptography. Chapter of the monography Boolean Methods and Models, Y. Crama and P. Hammer eds, Cambridge University Press, to appear (winter 2005-2006).\\[0pt]
[11] C. Carlet, P. Charpin and V. Zinoviev. Codes, bent functions and permutations suitable for DES-like cryptosystems. Designs, Codes and Cryptography, 15(2), pp. 125-156, 1998.\\[0pt]
[12] C. Carlet and C. Ding. Highly Nonlinear Mappings. Special Issue "Complexity Issues in Coding and Cryptography", dedicated to Prof. Harald Niederreiter on the occasion of his 60th birthday, Journal of Complexity 20, pp. 205-244, 2004.\\[0pt]
[13] F. Chabaud and S. Vaudenay. Links between differential and linear cryptanalysis, Advances in Cryptology -EUROCRYPT'94, Lecture Notes in Computer Science, Springer-Verlag, New York, 950, pp. 356-365, 1995.\\[0pt]
[14] N. Courtois and J. Pieprzyk. Cryptanalysis of block ciphers with overdefined systems of equations. Advances in cryptology-ASIACRYPT 2002, Lecture Notes in Computer Science 2501, pp. 267-287, Springer, 2003.\\[0pt]
[15] H. Dobbertin. One-to-One Highly Nonlinear Power Functions on $G F\left(2^{n}\right)$. Appl. Algebra Eng. Commun. Comput. 9 (2), pp. 139-152, 1998.\\[0pt]
[16] H. Dobbertin. Almost perfect nonlinear power functions over $G F\left(2^{n}\right)$ : the Niho case. Inform. and Comput., 151, pp. 57-72, 1999.\\[0pt]
[17] H. Dobbertin. Almost perfect nonlinear power functions over $G F\left(2^{n}\right)$ : the Welch case. IEEE Trans. Inform. Theory, 45, pp. 1271-1275, 1999.\\[0pt]
[18] H. Dobbertin. Almost perfect nonlinear power functions over $G F\left(2^{n}\right)$ : a new case for $n$ divisible by 5. D. Jungnickel and H. Niederreiter eds. Proceedings of Finite Fields and Applications FQ5, Augsburg, Germany, Springer, pp. 113-121, 2000.\\[0pt]
[19] R. Gold. Maximal recursive sequences with 3 -valued recursive crosscorrelation functions. IEEE Trans. Inform. Theory, 14, pp. 154-156, 1968.\\[0pt]
[20] H. Hollmann and Q. Xiang. A proof of the Welch and Niho conjectures on crosscorrelations of binary $m$-sequences. Finite Fields and Their Applications 7, pp. 253-286, 2001.\\[0pt]
[21] H. Janwa and R. Wilson. Hyperplane sections of Fermat varieties in $P^{3}$ in char. 2 and some applications to cyclic codes. Proceedings of AAECC-10, Lecture Notes in Computer Science, vol. 673, Berlin, Springer-Verlag, pp. 180-194, 1993.\\[0pt]
[22] T. Kasami. The weight enumerators for several classes of subcodes of the second order binary Reed-Muller codes. Inform. and Control, 18, pp. 369-394, 1971.\\[0pt]
[23] G. Lachaud and J. Wolfmann. The Weights of the Orthogonals of the Extended Quadratic Binary Goppa Codes. IEEE Trans. Inform. Theory, vol. 36, pp. 686-692, 1990.\\[0pt]
[24] M. Matsui. Linear cryptanalysis method for DES cipher. Advances in CryptologyEUROCRYPT'93, Lecture Notes in Computer Science, Springer-Verlag, pp. 386-397, 1994.\\[0pt]
[25] K. Nyberg. On the construction of highly nonlinear permutations. Advances in Cryptography, EUROCRYPT'92, Lecture Notes in Computer Science, Springer-Verlag, 658, pp. 92-98, 1993.\\[0pt]
[26] K. Nyberg. Differentially uniform mappings for cryptography, Advances in Cryptography, EUROCRYPT'93, Lecture Notes in Computer Science, Springer-Verlag, New York, 765, pp. 55-64, 1994.\\[0pt]
[27] K. Nyberg. S-boxes and Round Functions with Controllable Linearity and Differential Uniformity. Proceedings of Fast Software Encryption 1994, Lecture Notes in Computer Science 1008, pp. 111-130, 1995.\\[0pt]
[28] A. Pott. Nonlinear functions in Abelian groups and relative difference sets. Discrete Applied Math. 138, pp. 177-193, 2004.\\[0pt]
[29] V. Sidelnikov. On mutual correlation of sequences, Soviet Math. Dokl., 12(1971), pp. 197-201.\\[0pt]
[30] H. M. Trachtenberg. On the cross-correlation functions of maximal linear recurring sequences. PhD Thesis, University of Southern California, 1970.


\end{document}